\documentclass[lineno]{jfm}
\usepackage{graphicx}
\usepackage{multirow} % Paquete necesario para unir filas
\usepackage{caption}
\usepackage{subcaption}
\usepackage{booktabs}
\usepackage{float}
\usepackage{cancel}
\usepackage{ulem} 
%\usepackage{epstopdf,epsfig}
\usepackage{newtxtext}
\usepackage{newtxmath}
\usepackage{natbib}
\usepackage{comment}
\usepackage{hyperref}
\hypersetup{
    colorlinks = true,
    urlcolor   = blue,
    citecolor  = black,
}
\newtheorem{lemma}{Lemma}
\newtheorem{corollary}{Corollary}
\newtheorem{remark}{Remark}
% ---- FIX: \RomanNumeralCaps was declared with an argument count but no body,
%      so LaTeX silently swallowed the following \linenumbers as its definition
%      and line numbering never took effect. The macro was unused; removed.
\linenumbers

\newcommand{\orn}[1]{{\color{red} \textbf{Edit by Oscar:} #1}}
\newcommand{\coor}[1]{{\color{blue} #1}}
\newcommand{\ofore}[1]{{\color{yellow} #1}}
\newcommand{\cbp}[1]{{\color{blue} \textbf{Edit by Carlos:} #1}}
% ---- ROBUSTNESS: \cbp takes its content as a macro argument. This happens to
%      work for multi-paragraph text and display math, but it is fragile (any
%      \verb, footnote or catcode change inside would break it) and it makes the
%      final text blue. Long editorial blocks use the environment form below.
\newenvironment{cbpblock}%
  {\par\color{blue}\noindent\textbf{Edit by Carlos:}\ \ignorespaces}%
  {\par}
\newenvironment{ornblock}%
  {\par\color{red}\noindent\textbf{Edit by Oscar:}\ \ignorespaces}%
  {\par}


% {\MakeUppercase{\romannumeral #1}}
\title{An analytical closure for effective hydraulic resistance in Windkessel haemodynamics with generalised-Newtonian rheology}
%\title{An analytical closure for effective resistance in Windkessel haemodynamics with generalized-Newtonian rheology under physiological and pathological blood conditions}

%\title{A rheologically adaptive Windkessel model for haemodynamics in physiological and pathological states"}

%\title{A non-Newtonian Windkessel framework: Analytical formulation and application to pathological hemorheology}

%\title{The R-Windkessel: An analytical non-Newtonian framework for physiological and pathological reduced haemodynamics}



% ---- FIX: affiliation keys were inconsistent (\aff{1} never defined, \aff{2}
%      declared three times) and the CUP/DAMTP entries are template boilerplate.
%      Renumbered; please replace \aff{2} with the correct institution.
\author{Carlos Brito-Pacheco\aff{1}, Milenko Espinosa\aff{2},
  Ernesto Castillo\aff{2}, Scarlett Espinoza\aff{2}, Alfredo Parra\aff{2}, Luis Berr\aff{2},
 \and Oscar Ruz\aff{2}\corresp{\email{oscar.ruz@usach.cl}}}

\affiliation{
\aff{1}LMAP, CNRS, Université de Pau et des Pays de l’Adour, Pau, France.
\aff{2}Departamento de Ingeniería Mecánica, Universidad de Santiago de Chile,
Santiago, Chile.}

\begin{document}
\maketitle

\begin{abstract}
% ---- Previous version (written for an earlier version of the work) ----
% We derive an analytical closure for effective hydraulic resistance in Windkessel haemodynamics with generalised-Newtonian rheology. The closure is obtained from generalised Poiseuille-flow relations and is expressed as a multiplicative correction to the Windkessel resistance, thereby embedding constitutive rheology into a reduced resistive law. It recovers the Newtonian limit and accommodates both shear-thinning and shear-thickening responses in a unified resistance framework. The formulation is developed for purely viscous constitutive laws commonly used in blood rheology, namely the power-law, Carreau--Yasuda, Cross and Quemada models. Although the closure is general in structure, it is assessed here in a haemodynamic setting, where rheology-dependent resistance is especially relevant for lumped models and distal boundary representations. Physiological and pathological blood conditions, including anaemia, polycythaemia and diabetes, are incorporated through constitutive parameterisations spanning a broad range of haematocrit values. The resulting fitted rheological curves and parameter sets are reported in unified tabulations for direct use in haemorheological and reduced-modelling applications. The haemodynamic consequences of the proposed closure are investigated numerically by coupling it to an established zero-dimensional cardiac chamber model, exploring normotensive and hypertensive loading conditions and performing sensitivity analyses with respect to vessel radius, length and wall stiffness. We also provide an anchor validation of the resistance closure and examine how constitutive model choice influences effective resistance and reduced haemodynamic predictions at fixed haematocrit. The framework provides a consistent route to construct rheology-informed Windkessel models and physiologically informed distal boundary closures.
\cbp{WE WRITE THE ABSTRACT AT THE END.}

% ---- PROPOSED ABSTRACT (aligned with the reoriented argument; take or leave) ----
% The resistive element of a Windkessel model is the point at which blood
% viscosity enters a reduced haemodynamic description, yet it is conventionally
% built from Poiseuille arguments at a single constant viscosity. That constant
% is not a material property but a modelling choice, and the choice is
% consequential: for a healthy Carreau--Yasuda fit, evaluating the resistance at
% the high-shear plateau rather than at the operating shear rate changes it by
% about ten per cent, and a resistance calibrated at rest is in error by
% thirty-seven per cent at one-fortieth of the resting flow. We derive an
% analytical closure that removes this arbitrariness by carrying the
% generalised-Newtonian constitutive law into the quasi-steady pressure--flow
% relation itself, yielding a flow-dependent effective resistance with an exact
% derivative. The closure is instantiated for the power-law, Carreau--Yasuda,
% Cross and Quemada models and calibrated against shear viscometry for healthy,
% hypertensive, type~2 diabetic and hyperviscous blood. We then show that the
% constitutive dependence collapses onto a single quantity, the resting wall
% shear rate of a compartment, and that this quantity is fixed by calibre, path
% length and pressure share alone, independent of the flow. In a normally
% perfused coronary bed the identity places every compartment between $10^2$ and
% $10^3\,\mathrm{s}^{-1}$, so shear thinning contributes only a few per cent and
% a hyperviscous state is nearly indistinguishable from a healthy one. The
% closure is exercised in a coupled zero-dimensional left-ventricle /
% three-dimensional coronary model in which the venous drainage pressure is
% raised to emulate elevated central venous pressure. This displaces the venular
% bed towards the aggregation regime and separates the two constitutive states,
% demonstrating that the value of a rheology-dependent resistance lies not in the
% size of the correction at a nominal operating point but in its validity away
% from that point.
\end{abstract}

% Paper hierarchy:
% (1) Primary contribution: an analytical, rheology-informed closure for
%     effective hydraulic resistance, with an exact derivative, that removes the
%     arbitrary constant-viscosity choice from reduced models.
% (2) Secondary contribution: a unified constitutive-to-resistance procedure
%     applicable across several purely viscous rheological laws, together with a
%     criterion -- the resting wall shear rate identity -- that tells the reader
%     a priori whether the correction can matter in a given vascular bed.
% (3) Application layer: a single coupled 0D-LV/3D-coronary testbed. The 0D
%     ventricular model is a driver, not an object of study; the closure is
%     exercised at the coronary outlets, and the pathological contrast is driven
%     by the venous drainage pressure, which is the accessible route to the
%     low-shear regime in a normally perfused bed.
% Therefore, the paper should read first as a fluid-mechanics/reduced-modelling
% contribution, and only then as a haemodynamic application paper.
%
% IMPORTANT framing note: the paper must NOT claim that non-Newtonian rheology
% strongly alters coronary haemodynamics at the nominal operating point. It does
% not, and the reason is derived here rather than asserted. The claim is that a
% constant-viscosity resistance is valid only at its calibration point.

\noindent\textbf{MSC Codes}: 76Z05; 76A05; 76M12

\section{Introduction}
\label{sec:introduction}

% Cardiovascular disease is still the leading cause of
% deaths for both men and women worldwide. 
% Many risk factors have been identified and cur-
% rent therapeutic efforts have been centered on
% addressing these risk factors. However, as of
% today, the role that blood viscosity plays in this
% disease has not yet received its due attention.
% Viscosity is a fundamental property of any fluid.
% Its important role in both normal individuals and
% patients afflicted with cardiovascular disease has
% been underestimated. Past and current research
% has reported the benefits in addressing this impor-
% tant factor; however, mainstream medicine has
% not appreciated or fully accepted this important
% measurement.

\cbp{Cardiovascular disease constitutes one of the principal causes of mortality worldwide \citep{Roth2020GBD}. In the mechanical description of the circulation, the apparent viscosity of blood governs the relation between the pressure drop and the volumetric flow rate, and thereby contributes directly to the hydraulic resistance opposed to cardiac ejection; a variation in viscosity is, at fixed geometry, a variation in that resistance \citep{Chien1987,Mozos2015}. Blood, however, is not a Newtonian fluid -- its apparent viscosity depends on the local shear rate and, strongly and nonlinearly, on the haematocrit \citep{Chien1987,Nader2019}. The viscosity effectively realised in vivo is therefore not a fixed material constant, but varies over the cardiac cycle, through the shear-rate field, and between physiological and pathological states, through the haematocrit and the aggregation state of the erythrocytes. These dependencies motivate an explicit and quantitative treatment of blood viscosity in haemodynamic modelling, and more specifically in reduced descriptions, in which the flow resistance is conventionally evaluated at a single constant viscosity.}

% Blood is a concentrated suspension of deformable cells in an essentially Newtonian plasma, and its macroscopic behaviour is governed by suspension mechanics, including erythrocyte aggregation and disaggregation, deformation and alignment under shear.

\cbp{This dependence of the apparent viscosity on shear rate and haematocrit originates in the microstructure of blood, which is a concentrated suspension of deformable cells in an essentially Newtonian plasma \citep{FORMAGGIA-SPRINGER,Sequeria2018}; the macroscopic response is governed by suspension mechanics, in particular the aggregation and disaggregation, deformation and alignment of erythrocytes under shear.}

In haemodynamic modelling, the most widely adopted constitutive descriptions belong to the purely viscous generalised-Newtonian class, in which the deviatoric stress is proportional to the rate-of-deformation tensor and the apparent viscosity depends on shear rate and, directly or indirectly, on haematocrit \citep{FORMAGGIA-SPRINGER,Sequeria2018,Quemada1977,Popel2005}. Common choices include simple phenomenological relations, such as the power-law model, plateau-preserving models such as Cross and Carreau--Yasuda, and microstructure-motivated formulations such as Quemada \citep{Quemada1977,Popel2005,Neofytou2004,Mimouni2016}. Beyond purely viscous behaviour, experimental and modelling evidence indicates that blood may also exhibit yield-stress-like (viscoplastic) responses at low shear and viscoelastic effects associated with microstructural dynamics, particularly in regimes approaching stasis \citep{leuprecht2001computer}. The present work focuses on the purely viscous class as an analytically tractable and rheology-informed basis for resistance closures, yielding formulations that are both easy to use and systematically comparable across constitutive laws employed in haemodynamics.

Physiological and pathological conditions modify blood rheology through changes in the constitutive parameters of the chosen viscosity law. Haematocrit is a primary determinant, with viscosity increasing strongly and nonlinearly from severe anaemia to polycythaemia, with direct implications for cardiac workload and thrombotic risk \citep{Chien1987,Mozos2015,Nader2019,Isbister1994}. Anaemia reduces viscosity and can increase cardiac output \citep{murray1968circulatory,fowler1975blood}, whereas hyperviscosity syndromes, including polycythaemia-related states, are associated with impaired flow and thrombotic complications \citep{somer1993disorders,kwaan2003hyperviscosity,grenier2024red}. Diabetes mellitus has likewise been linked to haemorrheological alterations, including increased viscosity and enhanced non-Newtonian character in cohort measurements \citep{cho2008hemorheological,barnes1977blood,lowe1980blood,Sun2022}. Related metabolic factors, such as elevated cholesterol and triglycerides, have also been reported to promote aggregation and increase low-shear apparent viscosity \citep{moreno2015effect}. By contrast, other conditions act primarily through haemodynamics rather than intrinsic constitutive change. Hypertension chiefly alters pressure levels and pulsatility, while variations in vessel radius, length and wall stiffness modify flow partitioning and shear-rate fields. These drivers do not necessarily change the constitutive law itself, but they do shift the shear-rate range explored over the cardiac cycle and therefore alter the apparent viscosity effectively sampled by the flow.

This distinction has direct consequences for multiscale modelling. Introducing non-Newtonian physics in a three-dimensional domain is not only a matter of selecting a constitutive law in the governing equations, but also of prescribing outlet and distal boundary conditions informed by the same rheological mechanisms acting downstream. This is especially important because wall shear stress depends on both the near-wall velocity gradient and the local apparent viscosity, so rheological assumptions can influence shear-based biomarkers even when bulk flow quantities remain comparable. The issue therefore arises not only in fully lumped 0D descriptions, but also in one-dimensional wave-propagation models and in 0D--3D coupling strategies used to provide physiologically consistent distal boundary closures.

% Within reduced haemodynamics, Windkessel-type models and their network generalisations provide compact representations of distal vascular effects through lumped resistive and compliant elements \citep{frank1990basic,westerhof2009arterial}. These models are widely used both as standalone 0D descriptions and as outlet boundary closures for one-dimensional and three-dimensional solvers; extensions incorporating impedance concepts and wave-propagation effects have also been proposed and analysed from a fluid-mechanical standpoint \citep{bessems2007wave,mei2018fluid,westerhof2009arterial}. In such settings, resistive elements represent effective dissipation across unresolved vessel trees, while compliance captures storage due to wall distensibility. At the whole-heart level, Windkessel networks are commonly coupled to reduced cardiac chamber models, for example time-varying elastance-type descriptions, to generate closed-loop 0D dynamics suitable for parameter sweeps and sensitivity analyses. {\color{red}[Insert the corresponding heart 0D model reference here.]} This framework provides a natural setting in which to assess how constitutive rheology propagates into pressure--flow dynamics through effective resistance closures.

\cbp{Within reduced haemodynamics, Windkessel-type models and their network generalisations provide compact representations of distal vascular effects through lumped resistive and compliant elements \citep{frank1990basic,westerhof2009arterial}. Such models are employed both as standalone zero-dimensional descriptions and as outlet boundary conditions for one- and three-dimensional solvers, and extensions incorporating impedance and wave-propagation effects have been analysed from a fluid-mechanical standpoint \citep{bessems2007wave,mei2018fluid,westerhof2009arterial}. In these representations, the resistive element accounts for the effective dissipation across the unresolved vessel tree, and the compliant element for storage due to wall distensibility.

The resistive element is therefore the point at which the viscosity of blood enters a Windkessel description. In current practice, however, this resistance is constructed from Poiseuille-flow arguments at a constant viscosity, so that it depends only on the vessel geometry, even when the upstream model employs a generalised-Newtonian rheology, and even though the distal bed is the region in which non-Newtonian effects and haematocrit variability are most pronounced. The present work addresses this discrepancy.

We introduce a rheology-dependent resistance closure for Windkessel haemodynamics, derived from the quasi-steady, fully developed pressure–flow relation that underlies the classical constant-viscosity resistance; the additional element is the consistent incorporation of a generalised-Newtonian constitutive law within that same analytical setting. The closure is instantiated for four purely viscous laws, namely the power-law, Carreau–Yasuda, Cross and Quemada models, and is calibrated against shear-viscometry data for healthy, hypertensive, type 2 diabetic and hyperviscous blood over a range of haematocrit values \citep{Quemada1977,Popel2005,cho2008hemorheological,Sun2022}. The calibration is performed within a statistically informed procedure, motivated by the nonlinear dependence of viscosity on haematocrit and by the bias incurred when a constitutive law is fitted at a nominal mean value \citep{Chien1987,cho2008hemorheological,Westgard2024}.
}

It is worth stating at the outset what the closure does and does not change,
because the two are easily conflated. Replacing a constant viscosity by a
constitutive law does not, by itself, imply a large correction: whether it does
depends entirely on where the flow samples the viscosity curve. We show in
\S\ref{sec:RmuWindkessel} that for a compartment described by a calibre $R$, an
effective path length $L$ and a share $\Delta p$ of the pressure budget, the
resting wall shear rate is
\begin{equation}
\dot{\gamma}_{0}=\frac{R\,\Delta p}{2\,\mu_{\rm ref}\,L},
\label{eq:intro-gamma-identity}
\end{equation}
in which the flow rate cancels identically. The rheological operating point of
an unresolved bed is therefore fixed by morphometry and by the pressure
distribution, and is not a free modelling parameter. For a coronary
microvascular bed carrying a physiological pressure budget, \eqref{eq:intro-gamma-identity}
places every compartment---terminal arteriole, capillary, post-capillary venule
and collecting vein---between $10^{2}$ and $10^{3}\,\mathrm{s}^{-1}$, a regime in
which blood is close to Newtonian. This is confirmed independently by the
measured microvascular wall shear stress of $1$--$5\,\mathrm{Pa}$
\citep{Chien1987}, which upon division by the apparent viscosity returns the
same range.

Two consequences follow, and they organise the rest of the paper. First, the
correction at the nominal operating point is modest, of the order of a few per
cent above the high-shear plateau, and a hyperviscous constitutive state is
therefore only weakly distinguishable from a healthy one at rest. This is a
result rather than a limitation, and \eqref{eq:intro-gamma-identity} is the
instrument that establishes it a priori for any bed. Second, and this is where the
closure earns its place, a constant-viscosity resistance is accurate only at the
single flow rate at which it was calibrated, whereas the physiological range of
interest spans a factor of ten or more between low-flow states and maximal
vasodilation. Along that range the apparent viscosity of a healthy
Carreau--Yasuda fit varies by tens of per cent, so a resistance frozen at the
resting value carries an error that grows monotonically with distance from the
calibration point. The closure removes both this error and the arbitrariness of
the constant itself, which is not innocuous: for the same fit, evaluating the
resistance at the high-shear plateau instead of at the operating shear rate
already changes it by about ten per cent, an amount comparable to the shear-thinning
effect that motivated the exercise.

\cbp{The closure is exercised in a single computational setting, a coupled
zero-dimensional left-ventricle / three-dimensional coronary model in which the
reduced ventricular model of \citet{caruel2014dimensional} supplies the proximal
pressure and the intramyocardial pressure, and each coronary outlet carries an
$\mathcal{R}_{\mu}$-dependent reduced model. The pathological contrast is driven
by the venous drainage pressure: raising it, as occurs with elevated central
venous pressure in right-heart failure or tricuspid regurgitation, reduces the
venular driving gradient and displaces the venular bed towards the aggregation
regime, where the healthy and hyperviscous constitutive states separate. This
provides a physiologically grounded test of the closure that requires no change
of geometry and that probes precisely the shear range in which the constitutive
laws differ.}

\cbp{The remainder of the paper is organised as follows. Section~\ref{sec:models}
introduces the constitutive framework. Section~\ref{sec:Rmu} derives the
rheology-dependent resistance closure and specialises it to the four constitutive
laws considered. Section~\ref{sec:hemochar} describes the haemorheological
parameterisation and reports the calibration for the physiological and
pathological states. Section~\ref{sec:RmuWindkessel} introduces the
$\mathcal{R}_{\mu}$-Windkessel formulation, establishes the operating-point
identity~\eqref{eq:intro-gamma-identity} and quantifies the error incurred by a
constant-viscosity resistance away from its calibration point.
Section~\ref{sec:applications} describes the coupled 0D--3D coronary setting used
to exercise the closure and the venous-pressure protocol used to reach the
low-shear regime. Section~\ref{sec:results} reports the numerical results, and
Section~\ref{sec:conclusions} summarises the findings and discusses the
implications for rheology-consistent reduced modelling and for 0D--3D coupling.}

% In practice, however, Windkessel resistances are still commonly constructed from Poiseuille-flow arguments using a constant viscosity, yielding resistances that depend solely on geometry, even when the upstream model employs generalised-Newtonian rheology or when the distal bed is precisely where non-Newtonian effects and haematocrit variability are expected to be strongest. This is the central gap addressed in the present work. As in classical Windkessel resistance models, the present derivation relies on a quasi-steady fully developed relation between pressure drop and flow rate; the additional contribution here is the consistent incorporation of constitutive rheology within that same analytical setting. We therefore introduce a rheology-consistent, analytically tractable and rheology-informed resistance closure for Windkessel haemodynamics that can be parameterised systematically across physiological and pathological blood conditions. The closure takes the form of an analytical multiplicative factor that modifies the Windkessel resistance on the basis of generalised Poiseuille-flow relations, while the rheological calibration is rooted in standard shear-viscometry measurements. The framework is instantiated for several widely used purely viscous laws, namely the power-law, Carreau--Yasuda, Cross and Quemada models, and is reported through unified tabulations of fitted rheological curves and constitutive parameter sets spanning severe anaemia to polycythaemia and incorporating diabetes-related alterations \citep{Quemada1977,Popel2005,cho2008hemorheological,Sun2022}. Parameter estimation is performed in a statistically informed manner to improve robustness under cohort variability, motivated by the strong nonlinearity of viscosity with respect to haematocrit and the bias that can arise when fitting at a nominal mean \citep{Chien1987,cho2008hemorheological,Westgard2024}. The haemodynamic consequences of the proposed closure are assessed numerically by coupling it to an established zero-dimensional cardiac chamber model, exploring normotensive and hypertensive loading conditions and performing sensitivity analyses with respect to vessel radius, length and wall stiffness to represent arterial and arteriolar layers.

% The remainder of the paper is organised as follows. Section~\ref{sec:models} introduces the constitutive framework. Section~\ref{sec:Rmu} derives the rheology-dependent resistance closure. Section~\ref{sec:hemochar} details the haemorheological parameterisation strategies. Section~\ref{sec:RmuWindkessel} introduces the $\mathcal{R}_{\mu}$-Windkessel formulation. Section~\ref{sec:results} presents coupled 0D numerical experiments, sensitivity analyses and an anchor validation of the resistance closure, including an assessment of how constitutive model choice affects effective resistance and reduced haemodynamic predictions. Section~\ref{sec:conclusions} summarises the main findings and discusses implications for rheology-consistent reduced modelling and future 0D--3D coupling workflows.

\section{Mathematical models}
\label{sec:models}

\subsection{Navier--Stokes equations}

Let $\Omega\subset\mathbb{R}^{d}$, with $d=2$ or $3$, denote a bounded fluid domain with boundary $\Gamma=\partial\Omega$. We consider an incompressible flow described by the velocity field $\boldsymbol{u}:\Omega\times\mathbb{R}^{+}\rightarrow\mathbb{R}^{d}$ and the pressure field $p:\Omega\times\mathbb{R}^{+}\rightarrow\mathbb{R}$. The governing equations are
\begin{align}
\rho\,\partial_t\boldsymbol{u}+\rho\,\boldsymbol{u}\cdot\nabla\boldsymbol{u}-\nabla\cdot\boldsymbol{\sigma}(\boldsymbol{u},p)&=\boldsymbol{f}
\quad \text{in } \Omega,\label{eq:fluid-general-a}\\
\nabla\cdot\boldsymbol{u}&=0
\quad \text{in } \Omega,\label{eq:fluid-general-b}
\end{align}
where $\rho$ is the fluid density, $\boldsymbol{f}$ is the body force per unit volume, and $\boldsymbol{\sigma}$ is the Cauchy stress tensor. For a viscous incompressible fluid,
\begin{equation}
\boldsymbol{\sigma}=-p\boldsymbol{I}+\boldsymbol{\tau},\label{eq:cauchy-stress}
\end{equation}
where $\boldsymbol{I}$ is the identity tensor and $\boldsymbol{\tau}$ is the deviatoric stress tensor \citep{SPENCER-DOVER,Sequeria2018}. For a Newtonian fluid, the deviatoric stress is proportional to the rate-of-strain tensor
\begin{equation}
\boldsymbol{D}(\boldsymbol{u})=\frac{1}{2}\left(\nabla\boldsymbol{u}+(\nabla\boldsymbol{u})^{\mathrm{T}}\right),\label{eq:rate-strain}
\end{equation}
namely
\begin{equation}
\boldsymbol{\tau}=2\mu\,\boldsymbol{D}(\boldsymbol{u}),\label{eq:tau-newtonian}
\end{equation}
with $\mu$ the dynamic viscosity.

\subsection{Generalised-Newtonian constitutive laws}

Blood is a suspension of cellular elements in plasma. Since plasma is predominantly aqueous, it is commonly modelled as Newtonian, whereas the cellular phase induces non-Newtonian effects. Among the cellular components, erythrocytes dominate the rheological response because of their abundance and deformability. Their aggregation, disaggregation, deformation and alignment under shear give rise to the characteristic shear-thinning behaviour of blood \citep{haynes1959role,merrill1969rheology,chen2006non,ABUGATTAS2020529}. At low shear rates, rouleaux formation increases the apparent viscosity and may induce yield-stress-like behaviour; as shear increases, these aggregates break up and the apparent viscosity decreases \citep{FORMAGGIA-SPRINGER}.

% ---- Previous version ----
% To model this behaviour, we restrict attention to purely viscous generalised-Newtonian constitutive laws. In a general isotropic incompressible setting, the deviatoric stress may be expressed as in  Eq.\ref{eq:tau_general}; such materials are typically referred to as Reiner-Rivlin fluids \citep{FORMAGGIA-SPRINGER,Sequeria2018}
To model this behaviour, we restrict attention to purely viscous generalised-Newtonian constitutive laws. \cbp{In a general isotropic incompressible setting, the deviatoric stress of such materials, referred to as Reiner--Rivlin fluids \citep{FORMAGGIA-SPRINGER,Sequeria2018}, takes the form}
\begin{equation}
\boldsymbol{\tau}=\phi_1(I_1,I_2,I_3)\,\boldsymbol{D}(\boldsymbol{u})+\phi_2(I_1,I_2,I_3)\,\boldsymbol{D}(\boldsymbol{u})^2,\label{eq:tau_general}
\end{equation}
where $I_1$, $I_2$ and $I_3$ are the principal invariants of $\boldsymbol{D}(\boldsymbol{u})$,
\begin{equation}
I_1=\mathrm{tr}\bigl(\boldsymbol{D}(\boldsymbol{u})\bigr),\qquad
I_2=\frac{1}{2}\left[\mathrm{tr}\bigl(\boldsymbol{D}(\boldsymbol{u})\bigr)^2-\mathrm{tr}\bigl(\boldsymbol{D}(\boldsymbol{u})^2\bigr)\right],\qquad
I_3=\det\bigl(\boldsymbol{D}(\boldsymbol{u})\bigr).\label{eq:invariants}
\end{equation}
Because incompressibility implies $I_1=0$, and because the present work is concerned with purely viscous laws identified from viscometric measurements, we neglect second-order normal-stress contributions and retain the reduced form
\begin{equation}
\boldsymbol{\tau}=\phi_1(I_2)\,\boldsymbol{D}(\boldsymbol{u}).\label{eq:final_deviatoric}
\end{equation}
Introducing the shear rate
\begin{equation}
\dot{\gamma}=\sqrt{2\,\boldsymbol{D}(\boldsymbol{u}):\boldsymbol{D}(\boldsymbol{u})}=\sqrt{-4I_2},\label{eq:shear-rate}
\end{equation}
we obtain the standard generalised-Newtonian constitutive relation
\begin{equation}
\boldsymbol{\tau}=2\mu(\dot{\gamma})\,\boldsymbol{D}(\boldsymbol{u}),\label{eq:gnf}
\end{equation}
where $\mu$ is the apparent viscosity, or equivalently the shear-dependent viscosity.

% ---- Previous version ----
% The constitutive closure is therefore specified through the scalar function $\mu(\dot{\gamma})$. Particularly, in this work we consider four classical purely viscous models: Power-Law, Quemada, Cross, and Carreau-Yasuda, widely used in several applications, such as in hemodynamics \cite{}.
The constitutive closure is therefore specified through the scalar function $\mu(\dot{\gamma})$. \cbp{Specifically, in this work we consider four classical purely viscous models, namely the power-law, Cross, Carreau--Yasuda and Quemada models, widely used in haemodynamic applications~\citep{Quemada1977,Popel2005,Neofytou2004,Mimouni2016}.}

\subsection{Power-law model}

The power-law model is defined by
\begin{equation}
\mu=m\,\dot{\gamma}^{\,n-1},\label{eq:pl}
\end{equation}
where $m$ is the consistency index and $n$ is the flow index. It captures shear-thinning behaviour for $n<1$ and shear-thickening behaviour for $n>1$, but does not recover finite viscosity plateaux at very low or very high shear rates.

\subsection{Cross model}

The Cross model~\citep{Cross1965} is given by
\begin{equation}
\mu=\mu_{\infty}+\frac{\mu_{0}-\mu_{\infty}}{1+(\lambda\dot{\gamma})^{n}},\label{eq:cross}
\end{equation}
where $\mu_{0}$ and $\mu_{\infty}$ are the zero- and infinite-shear viscosities, $\lambda$ is a characteristic time scale, and $n$ controls the slope of the transition between the two limiting regimes. To avoid unnecessary symbol proliferation, the exponent $n$ is used in the power-law, Cross, and Carreau--Yasuda models, with its role defined locally in each constitutive expression.

\subsection{Carreau--Yasuda model}

The Carreau--Yasuda model~\citep{Carreau1972,Yasuda1981} reads
\begin{equation}
\mu=\mu_{\infty}+(\mu_{0}-\mu_{\infty})\left[1+(\lambda\dot{\gamma})^{a}\right]^{\frac{n-1}{a}},\label{eq:carreau-yasuda}
\end{equation}
where $a$ is the Yasuda parameter controlling the width of the transition region.

\subsection{Quemada model}

\begin{comment}
% ---- Previous version (varphi for volume fraction; $$-display; equation~\ref) ----
The Quemada model explicitly incorporates both shear rate and haematocrit, and is therefore especially attractive in haemorheological applications \citep{Quemada1977,Popel2005}. It is written as
\begin{equation}
\mu=\mu_{F}\left(1-\frac{1}{2}\frac{k_{0}+k_{\infty}\sqrt{\dot{\gamma}/\dot{\gamma}_{c}}}{1+\sqrt{\dot{\gamma}/\dot{\gamma}_{c}}}\,\varphi\right)^{-2},\label{eq:quemada}
\end{equation}
where $\mu_{F}$ is the viscosity of the suspending fluid, $\varphi$ is the dispersed-phase volume fraction, and $\dot{\gamma}_{c}$ is a critical shear rate separating aggregation- and dispersion-dominated regimes. The parameters $k_0$ and $k_{\infty}$ characterise the interaction state of the suspension at low and high shear rates, respectively.To align this formulation with the Cross and Carreau-Yasuda models, these states can be directly related to the asymptotic zero- and infinite-shear viscosities. Evaluating the limits $\mu_0 = \lim\limits_{\dot{\gamma} \to 0} \mu$ and $\mu_\infty = \lim\limits_{\dot{\gamma} \to \infty} \mu$ of equation~\ref{eq:quemada} yields
$$\mu_0 = \mu_F \left( 1 - \frac{1}{2} k_0 \varphi \right)^{-2} \quad \text{and} \quad \mu_\infty = \mu_F \left( 1 - \frac{1}{2} k_\infty \varphi \right)^{-2}.$$
\end{comment}

The Quemada model explicitly incorporates both shear rate and haematocrit, and is therefore \cbp{particularly suited to} haemorheological applications \citep{Quemada1977,Popel2005}. It is written as
\begin{equation}
\mu=\mu_{F}\left(1-\frac{1}{2}\frac{k_{0}+k_{\infty}\sqrt{\dot{\gamma}/\dot{\gamma}_{c}}}{1+\sqrt{\dot{\gamma}/\dot{\gamma}_{c}}}\,\varphi\right)^{-2},\label{eq:quemada}
\end{equation}
where $\mu_{F}$ is the viscosity of the suspending fluid, $\varphi$ is the dispersed-phase volume fraction, and $\dot{\gamma}_{c}$ is a critical shear rate separating aggregation- and dispersion-dominated regimes. The parameters $k_0$ and $k_{\infty}$ characterise the interaction state of the suspension at low and high shear rates, respectively.\cbp{ To align this formulation with the Cross and Carreau--Yasuda models, these states can be directly related to the asymptotic zero- and infinite-shear viscosities. Evaluating the limits $\mu_0 = \lim\limits_{\dot{\gamma} \to 0} \mu$ and $\mu_\infty = \lim\limits_{\dot{\gamma} \to \infty} \mu$ of \eqref{eq:quemada} yields}
\begin{equation}
\cbp{\mu_0 = \mu_F \left( 1 - \tfrac{1}{2} k_0 \varphi \right)^{-2} \quad \text{and} \quad \mu_\infty = \mu_F \left( 1 - \tfrac{1}{2} k_\infty \varphi \right)^{-2}.}\label{eq:quemada_limits}
\end{equation}

These constitutive laws provide the rheological input from which the effective resistance closure is derived in the following section.

\section{Generalised hydraulic resistance \texorpdfstring{$\mathcal{R}_{\mu}$}{Rmu}}
\label{sec:Rmu}

Consider steady, fully developed, laminar flow of an incompressible isothermal fluid through a rigid circular pipe of radius $R$ and length $L$. For a Newtonian fluid, the axial velocity profile is
\begin{equation}
v_z(r)=\frac{\Delta p}{4\mu L}(R^2-r^2),
\label{eq:v_newton}
\end{equation}
where $\Delta p$ is the pressure drop and $\mu$ is the dynamic viscosity. By contrast, the shear-stress distribution follows from the axial momentum balance alone and is independent of the constitutive law,
\begin{equation}
\tau(r)=\frac{r}{2}\frac{\Delta p}{L},
\label{eq:tau_dist}
\end{equation}
so that the wall shear stress is
\begin{equation}
\tau_w=\frac{R\,\Delta p}{2L}.
\label{eq:tau_w}
\end{equation}

The volumetric flow rate is defined by
\begin{equation}
Q=2\pi\int_{0}^{R}v_z(r)\,r\,\mathrm{d}r.
\label{eq:Q_def}
\end{equation}
For the Newtonian case, this yields the Hagen--Poiseuille relation,
\[
Q=\frac{\pi R^4}{8\mu L}\Delta p,
\]
which may be written as
\[
Q=\frac{\Delta p}{\mathcal{R}},
\qquad
\mathcal{R}=\frac{8\mu L}{\pi R^4}.
\]

For generalised-Newtonian fluids, explicit determination of $v_z(r)$ is not required. Instead, we use the Weissenberg--Rabinowitsch--Mooney--Schofield (WRMS) method \citep{Rabinowitsch1929,Bird1987}. Integrating \eqref{eq:Q_def} by parts and applying the no-slip condition gives
\begin{equation}
Q=\pi\int_{0}^{R}r^2\dot{\gamma}\,\mathrm{d}r,
\end{equation}
where $\dot{\gamma}$ is the shear rate. Using \eqref{eq:tau_dist} to change variables from $r$ to $\tau$ then yields
\begin{equation}
Q=\frac{\pi R^3}{\tau_w^3}\int_{0}^{\tau_w}\tau^2\dot{\gamma}(\tau)\,\mathrm{d}\tau.
\label{eq:WRMS_general}
\end{equation}
We therefore define a rheology-dependent hydraulic resistance $\mathcal{R}_{\mu}$ by
\begin{equation}
Q=\frac{\Delta p}{\mathcal{R}_{\mu}},
\label{eq:Q_Rmu}
\end{equation}
which, together with \eqref{eq:tau_w} and \eqref{eq:WRMS_general}, gives
\begin{equation}
\mathcal{R}_{\mu}
=
\frac{\Delta p\,\tau_w^3}{\pi R^3 I(\tau_w)}
=
\frac{\Delta p^4}{8\pi L^3 I(\tau_w)},
\label{eq:generalized-hydraulic-resistance}
\end{equation}
where
\begin{equation}
I(\tau_w)=\int_{0}^{\tau_w}\tau^2\dot{\gamma}(\tau)\,\mathrm{d}\tau
\label{eq:rheological_integral}
\end{equation}
is the rheological integral. The constitutive dependence of the effective resistance is therefore entirely contained in the mapping $\dot{\gamma}=\dot{\gamma}(\tau)$. We now propose the rheological integral for the viscous models presented before.

\subsection{Power-law hydraulic resistance}

The Power-Law fluid viscosity is given by,
\begin{equation}
\tau=m\dot{\gamma}^{n},
\qquad
\dot{\gamma}=\left(\frac{\tau}{m}\right)^{1/n},
\label{eq:pl_tau_gamma}
\end{equation}
whose substitution into \eqref{eq:rheological_integral} yields
\begin{equation}
I_{\mathrm{PL}}
=
\left(\frac{1}{m}\right)^{1/n}\frac{n}{3n+1}\tau_w^{(3n+1)/n},
\label{eq:I_PL}
\end{equation}
and hence
\begin{equation}
\mathcal{R}_{\mu}^{\mathrm{PL}}
=
\frac{\Delta p^4}{8\pi L^3 I_{\mathrm{PL}}(\tau_w)}.
\label{eq:R_PL}
\end{equation}

\subsection{Carreau--Yasuda hydraulic resistance}

The fluid viscosity given by the Carreau--Yasuda model is defined as,
\[
\mu(\dot{\gamma})
=
\mu_{\infty}
+
(\mu_0-\mu_{\infty})
\left[1+(\lambda\dot{\gamma})^a\right]^{(n-1)/a},
\]
where the relation $\tau=\mu(\dot{\gamma})\dot{\gamma}$ is implicit, so the rheological integral is expressed in terms of $\dot{\gamma}$. Letting $\delta=\mu_0-\mu_{\infty}$, one has
\begin{equation}
I_{\mathrm{CY}}
=
\int_{0}^{\dot{\gamma}_w}
\dot{\gamma}^3
\left[\mu_{\infty}+\delta\left(1+(\lambda\dot{\gamma})^a\right)^{\frac{n-1}{a}}\right]^2
\frac{\mathrm{d}\tau}{\mathrm{d}\dot{\gamma}}\,\mathrm{d}\dot{\gamma},
\label{eq:I_CY_start}
\end{equation}
with
% ---- Previous version (coefficient error: middle term should be \delta n, not \delta) ----
% \begin{equation}
% \frac{\mathrm{d}\tau}{\mathrm{d}\dot{\gamma}} = \mu_{\infty}
% +\delta \left(1+\lambda^a\dot{\gamma}^a\right)^{\frac{n-1}{a}}
% -\delta(n-1)\left(1+\lambda^a\dot{\gamma}^a\right)^{\frac{n-1-a}{a}}.
% \label{eq:dtau_dgamma_CY}
% \end{equation}
\begin{equation}
\frac{\mathrm{d}\tau}{\mathrm{d}\dot{\gamma}}
=
\mu_{\infty}
+\delta \cbp{n} \left(1+\lambda^a\dot{\gamma}^a\right)^{\frac{n-1}{a}}
-\delta(n-1)\left(1+\lambda^a\dot{\gamma}^a\right)^{\frac{n-1-a}{a}}.
\label{eq:dtau_dgamma_CY}
\end{equation}
\begin{comment}
% ---- Previous version (CY J-family without Cr/CY superscript) ----
Following \citet{Sochi2015}, this admits an exact representation in terms of Gauss hypergeometric functions. Defining
\begin{equation}
J_k
=
\int_{0}^{\dot{\gamma}_w}\dot{\gamma}^3(1+\lambda^a\dot{\gamma}^a)^k\,\mathrm{d}\dot{\gamma}
=
\frac{\dot{\gamma}_w^4}{4}\,
{}_2F_1\!\left(-k,\frac{4}{a};1+\frac{4}{a};-\lambda^a\dot{\gamma}_w^a\right),
\label{eq:Jk_CY}
\end{equation}
the integral becomes
\begin{align}
I_{\mathrm{CY}}
&=\mu_{\infty}^3J_0
+\mu_{\infty}^2\delta(n+2)J_{(n-1)/a}
-\mu_{\infty}^2\delta(n-1)J_{(n-1-a)/a}\nonumber\\
&\quad
+\mu_{\infty}\delta^2(2n+1)J_{(2n-2)/a}
-2\mu_{\infty}\delta^2(n-1)J_{(2n-2-a)/a}\nonumber\\
&\quad
+\delta^3nJ_{(3n-3)/a}
-\delta^3(n-1)J_{(3n-3-a)/a}.
\label{eq:I_CY_final}
\end{align}
\end{comment}

\cbp{Following \citet{Sochi2015}, this admits an exact representation in terms of Gauss hypergeometric functions. Defining
\begin{equation}
J_k^{\mathrm{CY}}
=
\int_{0}^{\dot{\gamma}_w}\dot{\gamma}^3(1+\lambda^a\dot{\gamma}^a)^k\,\mathrm{d}\dot{\gamma}
=
\frac{\dot{\gamma}_w^4}{4}\,
{}_2F_1\!\left(-k,\frac{4}{a};1+\frac{4}{a};-\lambda^a\dot{\gamma}_w^a\right),
\label{eq:Jk_CY}
\end{equation}
the integral becomes
\begin{align}
I_{\mathrm{CY}}
&=\mu_{\infty}^3J_0^{\mathrm{CY}}
+\mu_{\infty}^2\delta(n+2)J_{(n-1)/a}^{\mathrm{CY}}
-\mu_{\infty}^2\delta(n-1)J_{(n-1-a)/a}^{\mathrm{CY}}\nonumber\\
&\quad
+\mu_{\infty}\delta^2(2n+1)J_{(2n-2)/a}^{\mathrm{CY}}
-2\mu_{\infty}\delta^2(n-1)J_{(2n-2-a)/a}^{\mathrm{CY}}\nonumber\\
&\quad
+\delta^3nJ_{(3n-3)/a}^{\mathrm{CY}}
-\delta^3(n-1)J_{(3n-3-a)/a}^{\mathrm{CY}}.
\label{eq:I_CY_final}
\end{align}}
The associated hydraulic resistance is
\begin{equation}
\mathcal{R}_{\mu}^{\mathrm{CY}}
=
\frac{\Delta p^4}{8\pi L^3 I_{\mathrm{CY}}(\tau_w)}.
\label{eq:R_CY}
\end{equation}

\subsection{Cross hydraulic resistance}

\begin{comment}
% ---- Previous version (Cross exponent m; J-family index n) ----
For the Cross model, the fluid viscosity is represented by
\[
\mu(\dot{\gamma})
=
\mu_{\infty}
+
\frac{\delta}{1+(\lambda\dot{\gamma})^m},
\qquad
\delta=\mu_0-\mu_{\infty},
\]
where the differentiation of $\tau=\mu(\dot{\gamma})\dot{\gamma}$ gives
\begin{equation}
    d\tau = \left[ \mu_\infty + \frac{\delta}{1 + (\lambda \dot{\gamma})^m} - \frac{\delta m (\lambda \dot{\gamma})^m}{\left(1 + (\lambda \dot{\gamma})^m\right)^2} \right] d\dot{\gamma}\label{eq:dtau_dgamma_Cross}
\end{equation}
Substituting the preivous expression into the WRMS integral, yields
\begin{align}
    I = \int_0^{\dot{\gamma}_w} \dot{\gamma} & \left[ \mu_\infty^2 \dot{\gamma}^2 + \frac{2\mu_\infty \delta \dot{\gamma}^2}{1 + (\lambda \dot{\gamma})^m} + \frac{\delta^2 \dot{\gamma}^2}{\left(1 + (\lambda \dot{\gamma})^m\right)^2} \right] \cdot \nonumber \\
    & \left[ \mu_\infty + \frac{\delta}{1 + (\lambda \dot{\gamma})^m} - \frac{\delta m (\lambda \dot{\gamma})^m}{\left(1 + (\lambda \dot{\gamma})^m\right)^2} \right] d\dot{\gamma}
    \label{Int:Cross}
\end{align}

This integral can be expanded (see Section \ref{appA}) into several terms which together gives the following expression

\begin{align}
I &= \mu_{\infty}^3 J_0 + \left[ \mu_{\infty}^2 \delta (3 - m) \right] J_1 + \left[ \mu_{\infty} \delta^2 (3 - 2m) + \mu_{\infty}^2 m \delta \right] J_2 \nonumber \\
&+ \left[ \delta^3 (1 - m) + 2\mu_{\infty} \delta^2 m \right] J_3 + \left[ m \delta^3 \right] J_4
\end{align}
where $J_n$ corresponds to $J_n=\frac{\dot{\gamma}_w^4}{4} \, _2F_1\left(n, \frac{4}{m}; \frac{4}{m}+1; -(\lambda \dot{\gamma}_w)^m\right)$.

To ultimately derive $I_{Cr}$, \citet{Sochi2015} define the auxiliary variables $f=\lambda^m\dot{\gamma}_w^m$ and $g=1+f$.
\begin{equation}
f=\lambda^m\dot{\gamma}_w^m,
\qquad
g=1+f,
\label{eq:fg_cross}
\end{equation}
one may write
\begin{align}
I_{\mathrm{Cr}}
&=
\frac{\left( 2\delta^3 \zeta + \psi \right) \dot{\gamma}_w^4}{12m^2 g^3}
-\frac{\xi
{}_2F_1 \left(1, \frac{4}{m}; \frac{m+4}{m}; -f\right)\dot{\gamma}_w^4}{12m^2},
\label{eq:cross_integral_solution}
\end{align} where $\zeta=-m (2f^2 + 5f + 3) + 4g^2 + 2m^2$, $\psi=12m\delta^2 \mu_\infty g (m - g) + 3m^2\mu_\infty^2 g^2(4\delta + g)$, y $\xi=\delta^3 (m^2 - 6m + 8) + 3m\delta^2\mu_\infty (m - 4) + 3m^2\delta\mu_\infty^2$.
\end{comment}

\cbp{For the Cross model, the fluid viscosity is represented by
\[
\mu(\dot{\gamma})
=
\mu_{\infty}
+
\frac{\delta}{1+(\lambda\dot{\gamma})^n},
\qquad
\delta=\mu_0-\mu_{\infty},
\]
where the differentiation of $\tau=\mu(\dot{\gamma})\dot{\gamma}$ gives
\begin{equation}
    d\tau = \left[ \mu_\infty + \frac{\delta}{1 + (\lambda \dot{\gamma})^n} - \frac{\delta n (\lambda \dot{\gamma})^n}{\left(1 + (\lambda \dot{\gamma})^n\right)^2} \right] d\dot{\gamma}\label{eq:dtau_dgamma_Cross}
\end{equation}
Substituting the preivous expression into the WRMS integral, yields
\begin{align}
    I = \int_0^{\dot{\gamma}_w} \dot{\gamma} & \left[ \mu_\infty^2 \dot{\gamma}^2 + \frac{2\mu_\infty \delta \dot{\gamma}^2}{1 + (\lambda \dot{\gamma})^n} + \frac{\delta^2 \dot{\gamma}^2}{\left(1 + (\lambda \dot{\gamma})^n\right)^2} \right] \cdot \nonumber \\
    & \left[ \mu_\infty + \frac{\delta}{1 + (\lambda \dot{\gamma})^n} - \frac{\delta n (\lambda \dot{\gamma})^n}{\left(1 + (\lambda \dot{\gamma})^n\right)^2} \right] d\dot{\gamma}
    \label{Int:Cross}
\end{align}

This integral can be expanded (see Section \ref{appA}) into several terms which together gives the following expression

\begin{align}
I &= \mu_{\infty}^3 J_0^{\mathrm{Cr}} + \left[ \mu_{\infty}^2 \delta (3 - n) \right] J_1^{\mathrm{Cr}} + \left[ \mu_{\infty} \delta^2 (3 - 2n) + \mu_{\infty}^2 n \delta \right] J_2^{\mathrm{Cr}} \nonumber \\
&+ \left[ \delta^3 (1 - n) + 2\mu_{\infty} \delta^2 n \right] J_3^{\mathrm{Cr}} + \left[ n \delta^3 \right] J_4^{\mathrm{Cr}}
\end{align}
where $J_k^{\mathrm{Cr}}$ corresponds to $J_k^{\mathrm{Cr}}=\frac{\dot{\gamma}_w^4}{4} \, _2F_1\left(k, \frac{4}{n}; \frac{4}{n}+1; -(\lambda \dot{\gamma}_w)^n\right)$.

To ultimately derive $I_{Cr}$, \citet{Sochi2015} define the auxiliary variables $f=\lambda^n\dot{\gamma}_w^n$ and $g=1+f$.
\begin{equation}
f=\lambda^n\dot{\gamma}_w^n,
\qquad
g=1+f,
\label{eq:fg_cross}
\end{equation}
one may write
\begin{align}
I_{\mathrm{Cr}}
&=
\frac{\left( 2\delta^3 \zeta + \psi \right) \dot{\gamma}_w^4}{12n^2 g^3}
-\frac{\xi
{}_2F_1 \left(1, \frac{4}{n}; \frac{n+4}{n}; -f\right)\dot{\gamma}_w^4}{12n^2},
\label{eq:cross_integral_solution}
\end{align} where $\zeta=-n (2f^2 + 5f + 3) + 4g^2 + 2n^2$, $\psi=12n\delta^2 \mu_\infty g (n - g) + 3n^2\mu_\infty^2 g^2(4\delta + g)$, y $\xi=\delta^3 (n^2 - 6n + 8) + 3n\delta^2\mu_\infty (n - 4) + 3n^2\delta\mu_\infty^2$.} 

Finally, the hydraulic resistance for the cross model is defined as
\begin{equation}
\mathcal{R}_{\mu}^{\mathrm{Cr}}
=
\frac{\Delta p^4}{8\pi L^3 I_{\mathrm{Cr}}(\tau_w)}.
\label{eq:R_Cr}
\end{equation}

\subsection{Quemada hydraulic resistance}

For the Quemada model, the rheological integral was derived by \citet{Popel1993} using the WRMS method. In that formulation,
\begin{equation}
I_{\mathrm{Q}}
=
\left(\frac{\Delta p\,R}{2L}\right)^4
\frac{\left(1-\frac{1}{2}k\varphi\right)^2}{4\mu_F}
F(\alpha,q),
\label{eq:I_Q}
\end{equation}
where $k$ is a new parameter that encompasses the parameters originally defined in the Quemada model, $\frac{k_0+k_\infty \sqrt{\dot{\gamma}/\dot{\gamma_c}}}{1+\sqrt{\dot{\gamma}/\dot{\gamma_c}}}$, and $F(\alpha,q)$ is a dimensionless function involving polynomial coefficients from Quemada theory,
\begin{align}
F(\alpha,q)
&=
\frac{1}{2}\Bigg[
1-\frac{8}{7}\alpha(1+q)+\frac{4}{3}\alpha^2-\alpha^8P_7
+\left(1+\sum_{n=1}^{7}\alpha^nP_n\right)\sqrt{1-2\alpha q+\alpha^2}\nonumber\\
&\qquad\qquad
+\alpha^8P_8
\ln\left(
\frac{1-\alpha q+\sqrt{1-2\alpha q+\alpha^2}}{\alpha(1-q)}
\right)
\Bigg].
\label{eq:F_alpha_q}
\end{align}
Here, $\alpha$ and $q$ are dimensionless parameters defined in \citet{Popel1993}, and Polynomials $P_n$ are defined in Appendix \ref{appB}.
\begin{align*}
    \alpha = \frac{\sqrt{\tau_0} + \sqrt{\mu_\infty \lambda}}{\sqrt{\tau_w}}, \quad
    q = \frac{\sqrt{\tau_0} - \sqrt{\mu_\infty \lambda}}{\sqrt{\tau_0} + \sqrt{\mu_\infty \lambda}}
\end{align*}

The corresponding hydraulic resistance is
\begin{equation}
\mathcal{R}_{\mu}^{\mathrm{Q}}
=
\frac{\Delta p^4}{8\pi L^3 I_{\mathrm{Q}}(\tau_w)}.
\label{eq:R_Q}
\end{equation}

\section{Haemorheological characterisation}
\label{sec:hemochar}

\orn{The generalised hydraulic resistance $\mathcal{R}_{\mu}$ introduced in Section~\ref{sec:Rmu} provides a natural entry point for reduced-order descriptions of flows with non-constant viscosity. This is especially useful in haemodynamics, where complex vascular networks are routinely idealised as interconnected pipe segments represented through effective pressure--flow relations.}

\orn{Within this framework, blood is treated as a concentrated suspension of deformable red blood cells (RBCs) in a Newtonian plasma. Its macroscopic rheology is governed by the dynamics of the suspended phase. At low shear rates, plasma proteins, primarily fibrinogen, promote the formation of reversible stack-like aggregates known as rouleaux, which increase internal friction and apparent viscosity. As the shear stress increases, these aggregates disperse and the RBCs deform and align with the flow, producing the characteristic shear-thinning behaviour of blood.}


\subsection{Haematocrit influence and constitutive coupling}
\label{sec:haematocrit}

\orn{From a physiological standpoint, haematocrit is one of the principal parameters
governing red-cell interactions. A low haematocrit compromises oxygen transport
and drives the circulation towards a chronic hyperdynamic state which, if
sustained, may contribute to heart failure and increased cardiovascular
morbidity. A high haematocrit, by contrast, raises the blood viscosity strongly
and thereby the risk of ischaemic thrombotic events and, in extreme cases,
hyperviscosity syndrome with neurological and haemorrhagic manifestations
\citep{Mozos2015,Nader2019,Isbister1994}. The effective resistance is therefore
naturally regarded as a function not only of the pressure drop, vessel geometry
and shear-dependent viscosity, but also of the haematological state of the fluid,}
\begin{equation}
\mathcal{R}_{\mu}=\mathcal{R}_{\mu}(\Delta p,R,L,\mu,\varphi).
\label{eq:Rmu_phi}
\end{equation}

\begin{remark}
\orn{While the Quemada model \eqref{eq:quemada} is intrinsically formulated in terms
of the volume fraction~$\varphi$, the phenomenological models---power-law, Cross
and Carreau--Yasuda---require closure relations to couple their constitutive
parameters to the haematocrit. To ensure physical consistency across the shear
spectrum, the viscosity scales $(\mu_0,\mu_\infty,m)$ are made explicit functions
of~$\varphi$ using either the Krieger--Dougherty relation \citep{Krieger1959} or
an exponential empirical law \citep{CokeletA1963,CokeletB1963}.}
\end{remark}

\subsubsection{Coupling to haematocrit and population-consistent fitting}
\label{subsec:phi_coupling}

\orn{For the phenomenological models the constitutive scales must be made explicit
functions of~$\varphi$. We use the Krieger--Dougherty relation
\citep{Krieger1959} where a physically motivated packing correction suffices,
and, where a parameter shows stronger nonlinear dependence on haematocrit than a
single algebraic form can capture, an exponential representation}
\begin{equation}
\mu_i(\varphi)=c_1\exp\!\bigl(w(\varphi)\bigr),
\qquad
\mu_i(\varphi)=c_2+c_3\varphi,
\label{eq:phi_scaling}
\end{equation}
\orn{with $w(\varphi)$ a low-order polynomial and the linear form retained for weakly
$\varphi$-dependent parameters. The order of $w$ is fixed by calibration, so as
to preserve physical fidelity without introducing spurious degrees of freedom;
this follows the hybrid strategy of \citet{CokeletA1963,CokeletB1963}, who used
an exponential dependence for $\dot\gamma_c$ while keeping $k_0,k_\infty$ linear.
Here the exponential law is employed as an \emph{alternative} to
Krieger--Dougherty scaling, selected per parameter, rather than as an additional
fitted layer.}



\begin{remark}
\orn{In the microcirculation the haematocrit is not, in general, a conserved spatial
constant but exhibits a shear-dependent distribution associated with
shear-induced migration and the F\aa hr\ae us effect, so that $\varphi$ may vary
with the local strain rate, $\varphi=\varphi(\dot\gamma)$. In that regime the
coupling between the concentration field and the velocity gradient must be
resolved by evaluating the rheological integral of \S\ref{sec:Rmu} with a
shear-dependent volume fraction,
$I(\tau_w)=\int \tau^2\,\dot\gamma\bigl(\tau,\varphi(\dot\gamma)\bigr)\,
\mathrm{d}\tau$. The present closure retains a cross-sectionally
uniform~$\varphi$, which is adequate for the conduit- and distal-vessel scales
targeted here.}
\end{remark}
\begin{comment}
   Because haematocrit disperses appreciably across a clinical population, fitting a
constitutive law $\mathcal{M}(\dot\gamma,\varphi;\boldsymbol\theta)$ at the mean
haematocrit is biased whenever the viscosity is convex in~$\varphi$: by Jensen's
inequality $\mathbb{E}[\mu(\varphi)]\ge\mu(\mathbb{E}[\varphi])$, so the
mean-$\varphi$ evaluation systematically underestimates the population viscosity.
We therefore calibrate against the expectation of the constitutive law over the
haematocrit distribution,
\begin{equation}
\hat\mu(\dot\gamma;\boldsymbol\theta)
=\int_0^1 \mathcal{M}(\dot\gamma,\varphi;\boldsymbol\theta)\,P(\varphi)\,
\mathrm{d}\varphi,
\qquad
\boldsymbol\theta^{*}
=\arg\min_{\boldsymbol\theta}\sum_i
\bigl[\mu_{\mathrm{obs},i}-\hat\mu(\dot\gamma_i;\boldsymbol\theta)\bigr]^2,
\label{eq:expectation_fit}
\end{equation}
with $P(\varphi)\sim\mathcal{N}(\mu_\varphi,\sigma_\varphi)$ taken from the
reported ranges \citep{cho2008hemorheological}. This separates extrinsic haematocrit dispersion
from intrinsic constitutive behaviour, and the shaded bands of
figures~\ref{fig:fit_normal}--\ref{fig:fit_hiperv} are the image of
$P(\varphi)$ under each calibrated law. 
\end{comment}


\begin{comment}
    


\subsubsection{Krieger--Dougherty-based \texorpdfstring{$\mathcal{R}_{\mu}$}{Rmu} resistances}
\label{subsec:KD_Rmu}



\subsubsection{Exponential \texorpdfstring{$\mathcal{R}_{\mu}$}{Rmu} resistances}
\label{subsec:exp_Rmu}

Although Krieger--Dougherty scaling provides a physically motivated concentration correction, some constitutive parameters may exhibit stronger nonlinear dependence on haematocrit than can be captured by a single algebraic form. To account for this, we also consider a more flexible exponential representation,
\begin{equation}
\mu_i(\varphi)=c_1\exp\!\bigl(w(\varphi)\bigr),
\label{eq:exponential_gen}
\end{equation}
where $w(\varphi)$ is a polynomial function. This form provides a general framework for describing strong parameter sensitivity to haematocrit, especially near highly concentrated regimes. For parameters displaying weaker dependence on $\varphi$, a linear form may be sufficient,
\begin{equation}
\mu_i(\varphi)=c_2+c_3\varphi.
\label{eq:linear_gen}
\end{equation}
The order of $w(\varphi)$ is selected by calibration against experimental data, with the aim of preserving physical fidelity while avoiding unnecessary degrees of freedom.

A precedent for this hybrid strategy is given by \citet{CokeletA1963,CokeletB1963}, who used an exponential polynomial dependence for $\dot{\gamma}_c$ within a Quemada-type framework, while retaining linear relations for $k_0$ and $k_{\infty}$. In the present work, the same idea is embedded in the effective-resistance formulation. {\color{red}CHECK: state explicitly later whether this exponential strategy is used as an alternative to Krieger--Dougherty scaling or as an additional fitted parameterisation layer.}

\subsection{Probabilistic Distribution of \texorpdfstring{$\mathcal{R}_{\mu}$}{Rmu} resistances}
\label{subsec:fitted_Rmu}

In clinical populations, haematocrit exhibits significant inter-individual dispersion. Direct evaluation of constitutive models at the mean haematocrit may therefore introduce a systematic error associated with Jensen's inequality. In particular, if the viscosity function is convex with respect to $\varphi$, then
\begin{equation}
\mathbb{E}[\mu(\varphi)]\geq \mu(\mathbb{E}[\varphi]),
\label{eq:jensen}
\end{equation}
so that evaluating the viscosity at the mean haematocrit underestimates the expected viscosity of the population.

To address this issue, we define a general fitting framework for any constitutive model
\[
\mathcal{M}(\dot{\gamma},\varphi;\boldsymbol{\theta}),
\]
where $\boldsymbol{\theta}$ denotes the parameter vector to be estimated. The effective viscosity is defined as the expectation of the constitutive law with respect to the haematocrit distribution,
\begin{equation}
\hat{\mu}(\dot{\gamma};\boldsymbol{\theta})
=
\int_{0}^{1}\mathcal{M}(\dot{\gamma},\varphi;\boldsymbol{\theta})\,P(\varphi)\,\mathrm{d}\varphi,
\label{eq:mu_hat}
\end{equation}
where $P(\varphi)$ is the haematocrit probability density. In the present work, $P(\varphi)$ is modelled as a normal distribution consistent with clinical observations \citep{cho2008hemorheological},
\[
P(\varphi)\sim\mathcal{N}(\mu_{\varphi},\sigma_{\varphi}).
\]
The optimal parameter set is then obtained from the least-squares problem
\begin{equation}
\boldsymbol{\theta}^{*}
=
\arg\min_{\boldsymbol{\theta}}
\sum_i
\left[
\mu_{\mathrm{obs},i}
-
\hat{\mu}(\dot{\gamma}_i;\boldsymbol{\theta})
\right]^2.
\label{eq:theta_opt}
\end{equation}
This construction separates extrinsic haematocrit dispersion from intrinsic constitutive behaviour, leading to a more robust physiological characterisation and, consequently, a more reliable rheology-informed resistance.
\end{comment}

\orn{

\subsection{Blood pathologies and constitutive parameterisation}
\label{sec:pathologies}

The generalised resistance $\mathcal{R}_\mu$ of Section \ref{sec:Rmu} depends on
the constitutive law only through the shear-rate--shear-stress mapping
$\dot\gamma=\dot\gamma(\tau)$. A pathology therefore enters the reduced model not
as an additional term, but as a shift in the parameters of the chosen viscosity
law, and thus as a deformation of that mapping. In this section we calibrate the
four constitutive models against steady shear-viscometry data for four
haematological states---a healthy reference, arterial hypertension, type~2
diabetes mellitus and a hyperviscous state---each associated with a distinct
haematocrit range. The fitted parameters are collected in
table~\ref{tab:rheology}, and the corresponding fits are shown in
figures~\ref{fig:healthy}--\ref{fig:hyperviscosity}. Throughout, the
haematocrit~$\varphi$ is treated as the primary control parameter: it enters the
Quemada model intrinsically through the packing state of the suspension, whereas
for the power-law, Cross and Carreau--Yasuda models the viscosity scales
$(\mu_0,\mu_\infty,m)$ are coupled to $\varphi$ through either the
Krieger--Dougherty relation \cite{Krieger1959} or an exponential
empirical law \citep{CokeletA1963,CokeletB1963}.

Each dataset is fitted at its reported mean haematocrit, and the associated
physiological variability in~$\varphi$ is propagated to a shaded viscosity band in
figures~\ref{fig:healthy}--\ref{fig:hyperviscosity}. Because the
viscosity--haematocrit relation is strongly nonlinear, this band is asymmetric
and is not centred on the mean fit; reporting it explicitly guards against the
bias incurred when a single nominal $\varphi$ is used to represent a population
\cite{Chien1987,cho2008hemorheological,Westgard2024}.

\begin{remark}
In the microcirculation the haematocrit is not, in general, a conserved spatial
constant but exhibits a shear-dependent distribution associated with
shear-induced migration and the F\aa hr\ae us effect, so that $\varphi$ may vary
with the local strain rate, $\varphi=\varphi(\dot\gamma)$. In that regime the coupling
between the concentration field and the velocity gradient must be resolved by
evaluating the rheological integral of Section \ref{sec:Rmu} with a
shear-dependent volume fraction, $I(\dot\gamma_w)=\int \tau^2\,
\dot\gamma\bigl(\tau,\varphi(\dot\gamma)\bigr)\,\mathrm{d}\tau$. The present
closure retains a cross-sectionally uniform~$\varphi$, which is adequate for the
conduit- and distal-vessel scales targeted here.
\end{remark}

\subsubsection{Healthy reference}
Under physiological conditions blood rheology is governed by the high
deformability of the erythrocytes at a moderate haematocrit,
$0.40<\varphi<0.45$. This produces pronounced shear thinning, most marked in the
low-shear regime where reversible aggregation into rouleaux structures raises
the apparent viscosity. We take the healthy dataset of \cite{cho2008hemorheological}, at a
mean haematocrit $\varphi\approx0.42$, as the reference state. The apparent
viscosity decreases monotonically over $1$--$10^2\,\mathrm{s}^{-1}$
(figure~\ref{fig:healthy}), and all four models reproduce the mean trajectory
within the variability band ($\varphi=0.389$--$0.461$). The power-law fit, however,
remains a straight line in log--log coordinates and by construction cannot
recover the Newtonian plateaus at either shear-rate limit; the Cross,
Carreau--Yasuda and Quemada models capture the transition between the low- and
high-shear plateaus and provide the physically admissible baseline against which
the pathological states are compared.

\subsubsection{Arterial hypertension}
Hypertension-associated rheological changes are driven primarily by elevated
plasma fibrinogen, which promotes macromolecular bridging between erythrocytes;
combined with reduced cell deformability, this raises the bulk viscosity and
enhances aggregation. Using the data of \cite{Letcher1981} at $\varphi=0.444$
(figure~\ref{fig:hypertension}), the fits confirm strong shear thinning, with
the largest departure from Newtonian behaviour at low shear. Here the structural
deficiency of the power-law model is most visible: lacking both plateaus, it
deviates from the data at both shear-rate extremes. The Cross, Carreau--Yasuda
and Quemada models capture the nonlinear transition, although the Quemada fit
represents the low-shear dispersion less faithfully, a consequence of its
tendency to develop unphysical viscosity growth as the effective packing
fraction approaches unity (remark~\ref{rmk:quemada}).

\begin{remark}
\label{rmk:quemada}
The divergence of the Quemada viscosity as
\[
\tfrac{1}{2}\,
\frac{k_0+k_\infty\sqrt{\dot\gamma/\dot\gamma_c}}
{1+\sqrt{\dot\gamma/\dot\gamma_c}}\,\varphi \to 1
\]
is an artefact of the underlying minimum-energy-dissipation assumption and does
not reflect the physical response of erythrocyte suspensions at extreme
haematocrit. In practice this bounds the admissible parameter range of the
Quemada closure at the high-$\varphi$ end.
\end{remark}

\subsubsection{Type~2 diabetes mellitus}
The haemorheological signature of type~2 diabetes ($\varphi=0.410$--$0.490$) is a
marked elevation of bulk viscosity, most pronounced in the low-shear limit,
where reported values exceed those of healthy controls by more than $20\%$
\citep{cho2008hemorheological}. This is attributed to erythrocyte hyperaggregation and to a
reduction in cell deformability caused by chronic hyperglycaemia and the
glycation of membrane proteins, both of which raise the flow resistance of the
microvasculature \citep{Sun2022}. Fitting the data of \citet{cho2008hemorheological}
(figure~\ref{fig:diabetes}) yields a picture qualitatively similar to the
hypertensive case: the power-law model recovers the viscosity limits marginally
better than for hypertension but still fails to describe the intermediate-shear
transition, whereas the Cross, Carreau--Yasuda and Quemada models capture the
full nonlinear response.

\subsubsection{Hyperviscous state}
The hyperviscous state is parameterised from the data of \citet{Chien1966} at a
control haematocrit of $51.7\%$, where the non-Newtonian character of blood is
most strongly accentuated (figure~\ref{fig:hyperviscosity}). The apparent
viscosity rises steeply in the low-shear regime (down to
$\sim\!10^{-2}\,\mathrm{s}^{-1}$), producing the steepest shear-thinning
gradient of all four states. Under this severe nonlinearity the plateau-preserving
Cross, Carreau--Yasuda and Quemada models are clearly preferable to the power law.
Their behaviour nonetheless differs in an instructive way. The Quemada
uncertainty band is the most sensitive to the high volume fraction, broadening
into a dispersion that degrades its point-wise accuracy. The Carreau--Yasuda
model follows the experimental trajectory most closely, as expected from its
additional degree of freedom (the Yasuda exponent~$a$). The slightly larger
residual of the Cross model is not a failure but a consequence of its more
constrained functional form: with one parameter fewer, it imposes a stiffer
geometry on the transition region. This rigidity is desirable when the fitted
constants are to be interpreted structurally, since it reduces the risk of
overfitting to which the higher-dimensional Carreau--Yasuda form is more exposed.}

\begin{figure}
    \centering
    \begin{subfigure}[b]{0.48\textwidth}
        \includegraphics[width=\textwidth]{Images4/F1PL_Model_normal.pdf} 
        \caption{Power Law}
    \end{subfigure}
    \hfill
    \begin{subfigure}[b]{0.48\textwidth}
        \includegraphics[width=\textwidth]{Images4/F1Quemada_Model_normal.pdf} 
        \caption{Quemada Model}
    \end{subfigure}
    \medskip
    \begin{subfigure}[b]{0.48\textwidth}
        \includegraphics[width=\textwidth]{Images4/F1Cross_Model_normal.pdf} 
        \caption{Cross Model}
    \end{subfigure}
    \hfill
    \begin{subfigure}[b]{0.48\textwidth}
        \includegraphics[width=\textwidth]{Images4/F1CY_Model_normal.pdf} 
        \caption{Carreau-Yasuda Model}
    \end{subfigure}
    \caption{\textbf{Normal Subjects.} Reference rheological profiles for healthy controls. Data from \cite{cho2008hemorheological}.}
    \label{fig:healthy}
\end{figure} 




\begin{figure}
    \centering
    \begin{subfigure}[b]{0.48\textwidth}
        \includegraphics[width=\textwidth]{Images4/F2PL_Model_htn.pdf}
        \caption{Power Law}
    \end{subfigure}
    \hfill
    \begin{subfigure}[b]{0.48\textwidth}
        \includegraphics[width=\textwidth]{Images4/F2Quemada_Model_htn.pdf}
        \caption{Quemada Model}
    \end{subfigure}
    \medskip
    \begin{subfigure}[b]{0.48\textwidth}
        \includegraphics[width=\textwidth]{Images4/F2Cross_Model_htn.pdf}
        \caption{Cross Model}
    \end{subfigure}
    \hfill
    \begin{subfigure}[b]{0.48\textwidth}
        \includegraphics[width=\textwidth]{Images4/F2CY_Model_htn.pdf}
        \caption{Carreau-Yasuda Model}
    \end{subfigure}
    \caption{\textbf{Arterial Hypertension.} The complete multi-model profile confirms that cellular microrheology remains stable. Data from \cite{Letcher1981}.}
    \label{fig:hypertension}
\end{figure}



\begin{figure}
    \centering
    \begin{subfigure}[b]{0.48\textwidth}
        \includegraphics[width=\textwidth]{Images4/F3PL_Model_dm.pdf} 
        \caption{Power Law}
    \end{subfigure}
\hfill
    \begin{subfigure}[b]{0.48\textwidth}
        \includegraphics[width=\textwidth]{Images4/F3Quemada_Model_dm.pdf}
        \caption{Quemada Model}
    \end{subfigure}
    \medskip
    \begin{subfigure}[b]{0.48\textwidth}
        \includegraphics[width=\textwidth]{Images4/F3Cross_Model_dm.pdf}
        \caption{Cross Model}
    \end{subfigure}
    \hfill
    \begin{subfigure}[b]{0.48\textwidth}
        \includegraphics[width=\textwidth]{Images4/F3CY_Model_dm.pdf}
        \caption{Carreau-Yasuda Model}
    \end{subfigure}
    \caption{\textbf{Uncontrolled Diabetes.} Data from \cite{cho2008hemorheological}.}
    \label{fig:diabetes}
\end{figure}



\begin{figure}
    \centering
    \begin{subfigure}[b]{0.48\textwidth}
        \includegraphics[width=\textwidth]{Images4/F5PL_Model_Hyper.pdf}
        \caption{Power Law}
    \end{subfigure}
    \hfill
        \begin{subfigure}[b]{0.48\textwidth}
        \includegraphics[width=\textwidth]{Images4/F5Quemada_Model_Hyper.pdf}
        \caption{Quemada Model}
    \end{subfigure}    
    \medskip
    \begin{subfigure}[b]{0.48\textwidth}
        \includegraphics[width=\textwidth]{Images4/F5Cross_Model_Hyper.pdf}
        \caption{Cross Model}
    \end{subfigure}
    \hfill
    \begin{subfigure}[b]{0.48\textwidth}
        \includegraphics[width=\textwidth]{Images4/F5CY_Model_Hyper.pdf}
        \caption{Carreau-Yasuda Model}
    \end{subfigure}
    \caption{\textbf{hyperviscosity.} Data from \cite{Chien1966}}
    \label{fig:hyperviscosity}
\end{figure}






\section{The $\mathcal{R}_{\mu}$-Windkessel model}
\label{sec:RmuWindkessel}

Reduced-order lumped-parameter representations, commonly referred to as Windkessel models \citep{frank1990basic,olufsen2003deriving,westerhof2009arterial}, characterise the systemic arterial tree through a formal analogy with electro-mechanical circuits. In this framework, the complex spatio-temporal dynamics of the vasculature are mapped onto discrete elements representing hydraulic resistance, compliance, and inertance. The foundational two-element model comprises a linear resistance $\mathcal{R}$, accounting for viscous dissipation within the peripheral microvasculature, and a compliance $C$, representing the elastic energy storage and distensibility of the large conduit arteries. Higher-order extensions, such as the three-element model, incorporate an impedance term to account for proximal wave reflections, while four-element models introduce an inertance term to capture high-frequency pulsatile effects.

Irrespective of the specific Windkessel topology, the model may be generalised to incorporate the non-Newtonian behaviour of blood by replacing the constant (Poiseuille-type) resistance with a variable hydraulic resistance $\mathcal{R}_{\mu}$, as derived in \S\ref{sec:Rmu}. The $\mathcal{R}_{\mu}$-Windkessel model is thus governed by the non-linear constitutive flow relation
\begin{equation}
Q_{\mathrm{out}}=\frac{\Delta p}{\mathcal{R}_{\mu}(\dot{\gamma}, \varphi, \dots)},
\label{eq:WK_Rmu_outflow}
\end{equation}
where the resistance $\mathcal{R}_{\mu}$ is a function of the local shear rate $\dot{\gamma}$ and may depend on a vector of state-dependent parameters, including the haematocrit $\varphi$ and other pathology-specific rheological constants. This formulation provides a rigorous yet computationally efficient structure for investigating how microscopic rheological alterations propagate into macroscopic cardiovascular dynamics.

\subsection{The operating point of an unresolved bed}
\label{sec:operating-point}

Equation~\eqref{eq:WK_Rmu_outflow} is only as informative as the shear rate at
which it is evaluated, so it is worth establishing where an unresolved bed
actually sits on the viscosity curve. Represent a compartment as $N$ identical
vessels in parallel, of calibre $R$ and effective path length $L$, carrying a
branch flow $Q$ and dissipating a share $\Delta p$ of the pressure budget. At the
reference viscosity $\mu_{\rm ref}$ used for calibration, Poiseuille's relation
closes the multiplicity,
\begin{equation}
N=\frac{8\mu_{\rm ref}LQ}{\pi R^{4}\Delta p},
\label{eq:bed-multiplicity}
\end{equation}
and substituting into the definition of the wall shear rate,
$\dot{\gamma}_{0}=4Q/(N\pi R^{3})$, gives
\begin{equation}
\boxed{\;\dot{\gamma}_{0}=\frac{R\,\Delta p}{2\,\mu_{\rm ref}\,L}
=\frac{\tau_{w}}{\mu_{\rm ref}}\;}
\label{eq:gamma-identity}
\end{equation}
in which the flow rate has cancelled identically. The rheological operating point
of a compartment is therefore not a free parameter: it is fixed by the calibre,
the effective path length and the share of the pressure budget, and it is simply
the wall shear stress divided by the reference viscosity.

\begin{remark}
Equation~\eqref{eq:gamma-identity} is a useful screening tool. Since the
microvascular wall shear stress is an independently measured quantity of order
$1$--$5\,\mathrm{Pa}$, the operating shear rate of any bed carrying a
physiological pressure budget is of order $10^{2}$--$10^{3}\,\mathrm{s}^{-1}$. A
reduced model of such a bed will therefore sample the high-shear branch of any
purely viscous law, and the magnitude of the non-Newtonian correction can be
estimated before any simulation is performed.
\end{remark}

Because $\dot{\gamma}_{0}$ is fixed by morphometry, the pair
$(R,\bar{v})$---a calibre and a mean velocity, both directly accessible to
intravital microscopy---is the natural set of prescribed quantities, with
$\dot{\gamma}_{0}=4\bar{v}/R$, and with $N$, $L$ and the mean transit time
$T=L/\bar{v}$ following as predictions. Three measurable quantities
$(R,\bar{v},T)$ for two degrees of freedom leaves one consistency check: the
derived transit time must agree with the indicator-dilution value. Prescribing
$L$ instead and deriving $\dot{\gamma}_{0}$ would invert this hierarchy, since
$L$ is a lumped path through several generations in series and is not itself a
measurable quantity.

\subsection{Validity of a constant-viscosity resistance}
\label{sec:calibration-point}

The practical case for \eqref{eq:WK_Rmu_outflow} does not rest on the size of the
correction at a nominal operating point, but on the fact that a constant
resistance is exact only there. Writing the effective apparent viscosity as
$\mu_{\rm eff}(Q)=\mu_{\rm ap}(\dot{\gamma}_{0}Q/Q_{0})$, a resistance calibrated
at $Q_{0}$ incurs the relative error
\begin{equation}
e(Q)=\frac{\mu_{\rm ap}(\dot{\gamma}_{0})}{\mu_{\rm ap}(\dot{\gamma}_{0}Q/Q_{0})}-1,
\label{eq:calibration-error}
\end{equation}
which vanishes at $Q=Q_{0}$ by construction and grows monotonically away from it,
with opposite signs on either side. Two further remarks are in order. First, the
choice of the constant is itself consequential and rarely justified: for the
healthy fit of \S\ref{sec:hemochar}, evaluating the resistance at the high-shear
plateau $\mu_{\infty}$ rather than at $\mu_{\rm ap}(\dot{\gamma}_{0})$ already
changes it appreciably, by an amount comparable to the shear-thinning correction
itself. The closure removes this degree of freedom. Second, the closure supplies
the exact derivative
\begin{equation}
\frac{\partial Q}{\partial\Delta p}
=\frac{\pi R^{3}\dot{\gamma}_{w}-3Q}{\Delta p},
\label{eq:exact-derivative}
\end{equation}
obtained by differentiating \eqref{eq:WRMS_general} under the integral sign, so
that a Newton or monolithic coupling of the reduced model to a higher-dimensional
domain retains a consistent tangent. This is a practical advantage independent of
the magnitude of the rheological correction.

\begin{remark}
The closure is a statement about equilibrium rheology: $\mu_{\rm ap}$ is
evaluated as though the microstructure had fully adapted to the instantaneous
shear rate. Erythrocyte aggregation has a characteristic time of seconds, so the
low-shear branch should not be applied to transients shorter than that, such as a
flow reversal lasting tens of milliseconds. Where the results below approach
$\dot{\gamma}\to 0$ this restriction is stated explicitly.
\end{remark}

\section{Applications}
\label{sec:applications}

\begin{cbpblock}
The closure is exercised in a single computational setting: a coupled
LV-0D/coronary-3D model in which a three-dimensional coronary domain carries
rheology-dependent reduced models at its six terminal outlets. The reduced
left-ventricular model of \citet{caruel2014dimensional} appears in this setting
purely as a \emph{driver}: it supplies the proximal arterial pressure at the
coronary inlet and the intramyocardial pressure acting on the distal coronary
compartments. Its own systemic Windkessel stage is retained with
constant-viscosity resistances and is not an object of study here, so that the
$\mathcal{R}_{\mu}$ closure is exercised at one location only and the observed
sensitivity can be attributed unambiguously to the coronary outlet models. For
the same reason we do not report a normotensive/hypertensive sweep of the
systemic circuit, nor a sensitivity analysis of systemic vessel radius, length
and wall stiffness: those would probe the driver rather than the closure.

The coronary bed is chosen deliberately, and not because the correction is large
there. By the operating-point identity~\eqref{eq:gamma-identity} it is a bed in
which the correction is provably \emph{small} at rest, which makes it the
stringent case: any sensitivity recovered must come from displacing the operating
point rather than from a favourable choice of setting. The mechanism used to
displace it is the venous drainage pressure, described in
\S\ref{sec:pathology-protocol}.

All simulations were performed with the Rodin finite-element framework
\citep{Brito-Pacheco_Rodin_2023}. The implementations used here are those of the
\texttt{Heart/CCMLC2014} module and the
\texttt{examples/Heart/CoronaryArtery} coupled example.
\end{cbpblock}


\subsection{Reduced left-ventricular driver}
\label{sec:ApplicationI}

The proximal haemodynamic state of the coronary problem is supplied by a coupled
0D lumped-parameter model of the left heart and systemic circulation. Its role
here is to generate physiologically consistent waveforms for the arterial
pressure $p_{\rm ar}(t)$ and the ventricular pressure $p_{\rm LV}(t)$, which
enter the coronary model at the inlet and at the distal side of each outlet
respectively. It is described for completeness and reproducibility; its own
pressure--volume dynamics are not among the results reported in
\S\ref{sec:results}.

\subsubsection*{Ventricular and valvular dynamics}

The left ventricle is represented using the electromechanical framework of \cite{chapelle2012energy}, reduced under the assumptions of spherical symmetry and purely radial displacement \citep{caruel2014dimensional}. In contrast to traditional time-varying elastance models, this formulation explicitly accounts for the non-linear coupling between passive myocardial elasticity, active tension generation, and internal viscous damping. 

From a fluid-mechanical perspective, the flows through the mitral and aortic valves are treated as quasi-unidirectional and governed by pressure-driven gradients. The transitions between the filling, isovolumetric, and ejection phases are treated as state-dependent events. The conservation of mass within the ventricular cavity is governed by the kinematic relation
\begin{equation}
\frac{\mathrm{d}V}{\mathrm{d}t} = Q_{\rm in}(p_{\rm LA},p_{\rm LV}) - Q_{\rm out}(p_{\rm LV},p_{\rm ar}),
\end{equation}
where $p_{\rm LA}$ and $p_{\rm ar}$ denote the left-atrial and arterial pressures, respectively.

\subsubsection*{Systemic afterload}

The systemic afterload is modelled as a two-stage Windkessel circuit, partitioning the vasculature into proximal (conduit) and distal (microcirculatory) compartments. This arrangement is governed by the following system of coupled ordinary differential equations:
\begin{equation}
\begin{cases}
C_p \dfrac{\mathrm{d}p_{\rm ar}}{\mathrm{d}t} + \dfrac{p_{\rm ar} - p_d}{\mathcal{R}_{p}} = Q_{\rm out}, \\
C_d \dfrac{\mathrm{d}p_d}{\mathrm{d}t} + \dfrac{p_d - p_{\rm sv}}{\mathcal{R}_{d}} = \dfrac{p_{\rm ar} - p_d}{\mathcal{R}_{p}},
\end{cases}
\end{equation}
where $\{C_p, C_d\}$ and $\{\mathcal{R}_{p}, \mathcal{R}_{d}\}$ are the
stage-specific compliances and resistances. These resistances are kept at
constant viscosity on purpose. The closure of \S\ref{sec:Rmu} could equally be
applied here, and doing so would introduce a nonlinear feedback between cardiac
output and afterload; but it would also make the coronary response depend on two
simultaneously varying rheological elements, and the sensitivity reported in
\S\ref{sec:results} could no longer be attributed to the outlet closure alone.
The systemic stage is therefore held fixed, and the coronary inlet sees an
identical $p_{\rm ar}(t)$ in every case considered. This is verified in the
results as a control.

\subsection{Three-dimensional non-Newtonian coronary artery with
\texorpdfstring{$\mathcal{R}_{\mu}$}{Rmu}-Windkessel boundary conditions}
\label{sec:ApplicationII}

\cbp{
The second application considers a coupled LV-0D/coronary-3D model
(figure~\ref{fig:coronary-setup}). The
three-dimensional problem represents the flow in a coronary artery, while the
proximal haemodynamic state is supplied by the reduced left-ventricular model
introduced in \S\ref{sec:ApplicationI}. The coupling is one-way at the inlet:
the arterial pressure $p_{\rm ar}(t)$ produced by the LV-0D model is imposed
as the inlet pressure of the coronary domain. The ventricular pressure
$p_{\rm LV}(t)$ is also used in the coronary outlet models, where it plays
the role of an intramyocardial pressure acting on the distal coronary
compartments.

The reduced circulation associated with the LV model follows the pathway
\begin{equation}
p_{\rm LA}(t)
\xrightarrow{\,Q_{\rm LA}\,}
p_{\rm LV}
\xrightarrow{\,Q_{\rm ar}\,}
p_{\rm ar}
\xrightarrow{\,Q_p\,}
p_d
\xrightarrow{\,Q_d\,}
p_{\rm sv}.
\label{eq:lv0d-flow-topology}
\end{equation}
Here $p_{\rm LA}$ denotes the prescribed left-atrial pressure, $p_{\rm LV}$
the left-ventricular pressure, $p_{\rm ar}$ the proximal arterial pressure,
$p_d$ the systemic distal pressure and $p_{\rm sv}$ the systemic venous
pressure. The coronary model uses this same LV-0D state at two different
locations: $p_{\rm ar}$ drives the coronary inlet, whereas $p_{\rm LV}$
enters the distal side of each coronary outlet model. Thus, for each
terminal outlet $\Gamma_\mu^{(i)}$, $i=1,\ldots,6$, the coronary coupling can
be represented as
\begin{equation}
p_{\rm ar}
\longrightarrow
\Omega
\xrightarrow{\,Q_i\,}
p_{\rm out}^{(i)}
\xrightarrow{\,Q_{\mu,p}^{(i)}\,}
p_c^{(i)}
\xrightarrow{\,Q_{\mu,d}^{(i)}\,}
p_{\rm LV}.
\label{eq:coronary-flow-topology}
\end{equation}
The systemic distal pressure $p_d$ and the coronary capacitor pressures
$p_c^{(i)}$ are therefore distinct variables.

The computational domain $\Omega$ is a three-dimensional coronary artery
with one inlet, one wall boundary and six terminal outlets. The mesh is
rescaled to SI units and partitioned over 30  MPI ranks. It contains
35k tetrahedral elements,
11k vertices and
19k boundary faces. The spatial
discretisation uses Taylor--Hood finite elements~\citep{TaylorHood1973}, with continuous
piecewise-quadratic velocity and continuous piecewise-linear pressure,
\begin{equation}
\boldsymbol{u}_h\in [\mathbb{P}_2]^3,
\qquad
p_h\in \mathbb{P}_1.
\label{eq:coronary-taylor-hood}
\end{equation}
This gives 264768 velocity degrees of freedom,
14135 pressure degrees of freedom, and
\textbf{[TOTAL NUMBER OF DOFS]} total flow degrees of freedom.

The blood flow is modelled by the incompressible Navier--Stokes equations
with a Carreau--Yasuda generalised-Newtonian apparent viscosity. The
regularised shear rate is
\begin{equation}
\dot{\gamma}_{\varepsilon}
=
\left(
\varepsilon_{\gamma}^{2}
+
2\boldsymbol{D}(\boldsymbol{u}):\boldsymbol{D}(\boldsymbol{u})
\right)^{1/2},
\qquad
\boldsymbol{D}(\boldsymbol{u})
=
\frac{1}{2}
\left(
\nabla\boldsymbol{u}
+
\nabla\boldsymbol{u}^{T}
\right),
\label{eq:coronary-regularised-shear-rate}
\end{equation}
% ---- Previous version (restated the Carreau--Yasuda law of \eqref{eq:carreau-yasuda}) ----
% and the viscosity is evaluated as
% \begin{equation}
% \mu(\dot{\gamma}_{\varepsilon}) = \mu_{\infty} + (\mu_0-\mu_{\infty})\left[1+(\lambda\dot{\gamma}_{\varepsilon})^a\right]^{(n-1)/a}.
% \label{eq:coronary-cy-viscosity}
% \end{equation}
and the apparent viscosity follows the Carreau--Yasuda law~\eqref{eq:carreau-yasuda},
evaluated at the regularised shear rate $\dot{\gamma}_{\varepsilon}$.

The variational formulation includes a small pressure regularisation term
and inlet/outlet backflow stabilisation terms~\citep{EsmailyMoghadam2011}. The nonlinear convective term
and the viscosity-dependent terms are treated by an Oseen--Picard
linearisation~\citep{ElmanSilvesterWathen2014}. In practice, the advecting velocity, the Temam divergence
factor~\citep{Temam2001} and the Carreau--Yasuda viscosity are evaluated using the velocity
field from the previous time step. Each accepted time step therefore requires
the solution of one linearised three-dimensional saddle-point problem.

Time integration is first order. At each time step, the LV-0D model is first
advanced to obtain $p_{\rm LV}^{n+1}$ and $p_{\rm ar}^{n+1}$. The pressure
$p_{\rm ar}^{n+1}$ is then imposed at the coronary inlet and the
three-dimensional flow problem is solved. The resulting outlet fluxes
$Q_i^{n+1}$ are used to update the six outlet models, from which the outlet
pressures $p_{\rm out}^{(i),n+1}$ are reconstructed and imposed as traction
pressures on the corresponding terminal boundaries. The average
computational cost was 12 seconds per accepted
time step using 30 MPI ranks on
\textbf{[HARDWARE / CPU DESCRIPTION]}. The total simulated physical time was
2.55 seconds with time step
$\Delta t=0.001$ seconds.
}
% ---- NOTE: the outlet formulation below is deliberately NOT wrapped in \cbp{}.
%      It is final text rather than an editorial annotation, and it should not be
%      typeset in blue.

Each terminal coronary outlet carries an independent nonlinear
$\mathcal{R}_{\mu}$ reduced model with a single state, the microvascular
transmural pressure
\begin{equation}
p_{tm}^{(i)}=p_c^{(i)}-p_{\rm im},
\qquad
p_{\rm im}=\alpha\,p_{\rm LV},
\label{eq:coronary-transmural}
\end{equation}
where $p_c^{(i)}$ is the microvascular pressure and $\alpha$ the
volume-averaged transmission of cavity pressure to the myocardial tissue. The
sign convention is such that $Q_i>0$ denotes flow leaving the
three-dimensional domain, and $Q_{d}^{(i)}>0$ drainage towards the venous
reservoir at pressure $P_{\rm RA}$.

The compartment is intramyocardial, so the tissue pressure is not a source added
to the balance but the reference of the whole compartment. Its venous limb is
accordingly a Starling resistor whose throat closes at $p_{\rm im}$
\citep{Downey1975,Permutt1963}, giving
\begin{align}
C^{(i)}\frac{\mathrm{d}p_{tm}^{(i)}}{\mathrm{d}t}
&=Q_i-Q_{d}^{(i)},
\label{eq:coronary-node-balance}\\
Q_{d}^{(i)}
&=Q_{\mu,d}^{(i)}\!\left(p_{tm}^{(i)}+p_{\rm im}-p_{w}\right),
\qquad
p_{w}=P_{\rm RA}+\beta\left(p_{\rm im}-P_{\rm RA}\right)^{+},
\label{eq:coronary-starling}\\
p_{\rm out}^{(i)}
&=p_{\rm im}+p_{tm}^{(i)}+\mathcal{R}_{\mu,p}^{(i)}\!\left(Q_i\right)Q_i,
\label{eq:coronary-outlet-pressure}
\end{align}
in which $\beta\in[0,1]$ is the fraction of tissue pressure reaching the venous
throat. The limits are informative: $\beta=1$ is the strong waterfall, in which
$p_{\rm im}$ cancels from the venous driving pressure entirely, whereas $\beta=0$
recovers drainage against a fixed reference. The intermediate value reflects the
fact that the venular exit lies closer to the epicardium than the volume-averaged
tissue pressure of the bed, so it is exposed to a fraction of it. The strong
waterfall is not the modern consensus---capacitive mechanisms account for much of
the observed zero-flow behaviour \citep{Spaan1981,Klocke1985,Westerhof2006}---and
$\beta$ interpolates continuously between the two descriptions.

Two properties of \eqref{eq:coronary-node-balance}--\eqref{eq:coronary-outlet-pressure}
are worth noting, because they are what reduces the outlet to one state. First,
$p_{\rm im}$ appears explicitly in the outlet pressure
\eqref{eq:coronary-outlet-pressure}, so the systolic impediment to coronary
inflow requires no separate calibre modulation and no explicit pump source.
Second, in the waterfall regime $p_{\rm im}$ cancels from
\eqref{eq:coronary-starling}, so $p_{tm}$ becomes a relaxation variable with
time constant $C^{(i)}\mathcal{R}_{\mu,d}^{(i)}$, bounded below by zero and above
by $\mathcal{R}_{\mu,d}^{(i)}\max Q_i$. The stored volume $C^{(i)}p_{tm}^{(i)}$ is
therefore non-negative by construction, without a collapsible-tube law or a
non-negativity penalty.

Both branch flows are evaluated with the $\mathcal{R}_{\mu}$ closure of
\S\ref{sec:Rmu}, and the operating shear rates follow from
\eqref{eq:gamma-identity} rather than being prescribed. The proximal resistance is
assembled implicitly on the outlet boundary. Writing the resistive contribution to
the Neumann traction for a flat profile,
\begin{equation}
\int_{\Gamma_\mu^{(i)}}\!\mathcal{R}_{\mu,p}^{(i)}Q_i\,(\boldsymbol{v}\cdot\boldsymbol{n})
\,\mathrm{d}\Gamma
\;\simeq\;
\mathcal{R}_{\mu,p}^{(i)}A_i
\int_{\Gamma_\mu^{(i)}}\!(\boldsymbol{u}\cdot\boldsymbol{n})
(\boldsymbol{v}\cdot\boldsymbol{n})\,\mathrm{d}\Gamma ,
\label{eq:coronary-implicit-resistance}
\end{equation}
with $A_i$ the outlet area, gives a symmetric positive-semidefinite operator that
is exact for a flat profile. Assembling it rather than lagging it by one step
changes the amplification factor of the lagged coupling from
$|1-\mathcal{R}\Delta t/L_{3D}|$, which exceeds unity for the dominant distal
resistance, to $1/(1+\mathcal{R}\Delta t/L_{3D})<1$ for every $\Delta t$ and every
$\mathcal{R}$, where $L_{3D}$ is the inertance of the three-dimensional domain.
The coupling is then unconditionally stable and requires no outlet capacitor; this
matters because such a capacitor forms an undamped resonator with $L_{3D}$ and
contaminates the outlet flux with a spurious systolic component.

% ---- Previous version (re-derived the Section 3 WRMS law; superseded by a reference to \S\ref{sec:Rmu}) ----
% Both the proximal and distal outlet flows are evaluated with the same
% Carreau--Yasuda WRMS-type pressure--flow law. For a pressure drop
% $\Delta p$, an effective vessel radius $R$ and length $L$, the wall shear
% stress is
% \begin{equation}
% \tau_w = \frac{R|\Delta p|}{2L}, \label{eq:coronary-wall-shear-stress}
% \end{equation}
% and the wall shear rate is obtained from
% \begin{equation}
% \tau_w = \mu(\dot{\gamma}_w)\dot{\gamma}_w. \label{eq:coronary-wall-shear-rate}
% \end{equation}
% Assuming that the map $\dot{\gamma}\mapsto \mu(\dot{\gamma})\dot{\gamma}$ is monotone over the
% considered shear-rate interval, the corresponding flow is computed as
% \begin{equation}
% Q_{\mu}(\Delta p;R,L) = \operatorname{sign}(\Delta p) \frac{\pi R^3}{\tau_w^3}
% \int_0^{\dot{\gamma}_w} \dot{\gamma}^{3} \mu(\dot{\gamma})^2
% \frac{\mathrm{d}}{\mathrm{d}\dot{\gamma}}\left[ \mu(\dot{\gamma})\dot{\gamma} \right]\,\mathrm{d}\dot{\gamma}.
% \label{eq:coronary-wrms-flow-law}
% \end{equation}
For a pressure drop $\Delta p$ across an outlet segment of effective radius $R$
and length~$L$, the wall shear stress follows from \eqref{eq:tau_w} and the flow
from the rheological integral \eqref{eq:WRMS_general}, generalised here to signed
pressure drops,
\begin{equation}
Q_{\mu}(\Delta p;R,L)
=
\operatorname{sign}(\Delta p)\,
\frac{\pi R^3}{\tau_w^3}\,I_{\mathrm{CY}}(\dot{\gamma}_w),
\qquad
\tau_w=\frac{R|\Delta p|}{2L},
\label{eq:coronary-wrms-flow-law}
\end{equation}
where $I_{\mathrm{CY}}$ is the Carreau--Yasuda rheological integral
\eqref{eq:I_CY_final} and the wall shear rate $\dot{\gamma}_w$ solves
$\tau_w=\mu(\dot{\gamma}_w)\dot{\gamma}_w$.

Rather than evaluating \eqref{eq:coronary-wrms-flow-law} at every call, we exploit
a structural property of the closure. Writing the flow of a
generalised-Newtonian fluid as Hagen--Poiseuille with an apparent viscosity and
comparing with \eqref{eq:coronary-wrms-flow-law} gives
\begin{equation}
\mu_{\rm ap}(\tau_w)=\frac{\tau_w^{4}}{4\,I(\tau_w)},
\qquad
\dot{\gamma}_{\rm nom}=\frac{4Q}{\pi R^{3}}=\frac{\tau_w}{\mu_{\rm ap}},
\label{eq:mu-apparent}
\end{equation}
so that $\mu_{\rm ap}$ is a \emph{universal} function of the wall shear stress for
a given rheology: it carries no dependence on $R$, on $L$ or on the branch. It is
therefore tabulated once in $\dot{\gamma}_{\rm nom}$ and the reduced element only
interpolates. With $241$ logarithmically spaced nodes over
$\tau_w\in[10^{-6},10^{4}]\,\mathrm{Pa}$ the maximum log--log interpolation error
is $0.09\%$, against a per-call root solve for $\dot{\gamma}_w$ and a quadrature
of the rheological integral. This eliminates the nonlinear inner solves entirely,
and with them the need for any fallback approximation: the outlet update reduces to
a scalar Newton iteration on $p_{tm}^{(i)}$ whose residual derivative,
$C^{(i)}/\Delta t+\partial Q_{d}^{(i)}/\partial p_{tm}^{(i)}$, is strictly
positive, so the iteration converges globally from any iterate.

\subsubsection{Pathological protocol: venous drainage pressure}
\label{sec:pathology-protocol}

\cbp{
The constitutive states of \S\ref{sec:hemochar} differ chiefly at low shear, while
by \eqref{eq:gamma-identity} a normally perfused coronary bed operates two to
three decades above that range. Distinguishing them therefore requires displacing
the operating point, and the cheapest physiologically grounded way to do so is
through the venous drainage pressure $P_{\rm RA}$.

Elevated central venous pressure---in right-heart failure, tricuspid
regurgitation or constrictive pericarditis---raises the downstream reference of
the coronary venous drainage without altering any vessel geometry. Through
\eqref{eq:coronary-starling} it reduces the venular driving gradient, hence the
venular flow, hence by \eqref{eq:gamma-identity} the venular shear rate. It is
also the only accessible mechanism with this property: an increase in
microvascular resistance by capillary rarefaction leaves the per-vessel flow, and
therefore the shear rate, unchanged, since the flow and the number of parallel
vessels fall together; vasoconstriction reduces the shear rate only as the square
root of the resistance increase, because the calibre falls with the flow.

We therefore sweep $P_{\rm RA}$ across its clinical range at fixed geometry and
fixed systemic driver, for each of the four constitutive states of
\S\ref{sec:hemochar}, and report the resulting flow, waveform and shear-rate
statistics. This is a test of the closure rather than a search for a favourable
configuration: the sweep is monotone in a single physiological parameter, and the
prediction to be verified is that the constitutive states separate progressively
as the venular operating point descends through the aggregation transition.
Where the sweep drives $\dot{\gamma}$ below a few tens of $\mathrm{s}^{-1}$ the
equilibrium restriction noted in \S\ref{sec:calibration-point} applies and is
reported alongside the results.
}

\begin{figure}
    \centering
    \includegraphics[width=0.9\linewidth]{CoronarySetup.pdf}
    \caption{
    Schematic of the coupled LV-0D/coronary-3D model. The reduced LV model
    follows the pathway
    $p_{\rm LA}(t)\rightarrow p_{\rm LV}\rightarrow p_{\rm ar}
    \rightarrow p_d\rightarrow p_{\rm sv}$ and acts as a driver only: the
    arterial pressure $p_{\rm ar}$ is imposed at the inlet of the
    three-dimensional coronary domain, and the ventricular pressure
    $p_{\rm LV}$ sets the intramyocardial pressure
    $p_{\rm im}=\alpha p_{\rm LV}$ of the coronary outlet models. Each terminal
    outlet $\Gamma_\mu^{(i)}$ carries a single-state $\mathcal{R}_{\mu}$ model:
    a proximal rheology-dependent resistance $\mathcal{R}_{\mu,p}^{(i)}$
    assembled implicitly on the boundary, a microvascular compartment of
    compliance $C^{(i)}$ referenced to $p_{\rm im}$ with state
    $p_{tm}^{(i)}=p_c^{(i)}-p_{\rm im}$, and a distal rheology-dependent
    resistance $\mathcal{R}_{\mu,d}^{(i)}$ terminated by a Starling throat that
    closes at $p_{w}=P_{\rm RA}+\beta(p_{\rm im}-P_{\rm RA})^{+}$. The
    pathological protocol of \S\ref{sec:pathology-protocol} varies $P_{\rm RA}$
    at fixed geometry.
    \textbf{[FIGURE TO BE REDRAWN: the current \texttt{CoronarySetup.pdf} shows
    the superseded capacitor-plus-source outlet; it must be replaced by the
    single-state topology above.]}
    }
    \label{fig:coronary-setup}
\end{figure}
\section{Results}\label{sec:results}

\cbp{OUTLINE ONLY --- the numbered items below are the results to be produced,
in the order in which they build the argument. Indicative magnitudes in square
brackets come from the Carreau--Yasuda fits of table~\ref{tab:parametros_reologicos}
evaluated through the closure, and are placeholders to be replaced by the
simulation output.}

\begin{enumerate}

\item \textbf{Verification of the closure.} Newtonian limit
($\mu_0\!\to\!\mu_\infty$) recovered to machine precision; low- and high-shear
plateaux of $\mu_{\rm ap}$ recovered as $\mu_0$ and $\mu_\infty$; agreement
between the tabulated and directly integrated closure
[max.\ error $0.09\%$, mean $0.011\%$ over $241$ nodes]. One table, no figure.

\item \textbf{The operating point is not free.} Evaluation of
\eqref{eq:gamma-identity} for the four compartments of a coronary bed
(terminal arteriole, capillary, post-capillary venule, collecting vein) with
morphometric $R$ and the measured pressure budget, tabulated against the
independent estimate $\dot{\gamma}=\tau_w/\mu$ from
$\tau_w=1$--$5\,\mathrm{Pa}$. Expected outcome: all compartments in
$10^{2}$--$10^{3}\,\mathrm{s}^{-1}$. This is the screening result and it should
appear before any simulation. \emph{Include the transit-time consistency check
$T=L/\bar{v}$ against indicator dilution, and report it even where it fails}
[$T_a\!\approx\!8\,$s for $R=25\,\mu$m, $\bar v=5\,$mm/s, against $1$--$2\,$s
measured]: it bounds the validity of representing a branching tree by a single
effective segment, and is more informative stated than omitted.

\item \textbf{Constitutive contrast at the operating point.} $\mu_{\rm ap}$ for
the four states of table~\ref{tab:parametros_reologicos} across
$\dot{\gamma}\in[1,10^{3}]\,\mathrm{s}^{-1}$, with the ratio to the healthy state.
Expected outcome: hyperviscous/healthy
[$\approx\!1.17$ at $400\,\mathrm{s}^{-1}$, $\approx\!1.52$ at
$10\,\mathrm{s}^{-1}$]; diabetes and hypertension cross the healthy curve.
One figure, log--log, four curves plus a shaded band marking the coronary
operating range. This figure states the difficulty and motivates item~5.

\item \textbf{Error of a constant-viscosity resistance.} $e(Q)$ from
\eqref{eq:calibration-error} over $Q/Q_0\in[1/40,4]$ for each constitutive state,
calibrated at rest. Expected outcome for the healthy fit
[$-37\%$ at $Q/Q_0=1/40$, $-12.5\%$ at $1/8$, $+2.3\%$ at $4$]. Add the
plateau-versus-operating-point comparison ($\mu_\infty$ against
$\mu_{\rm ap}(\dot{\gamma}_0)$) as a horizontal reference line. \emph{This is the
principal methodological result of the paper} and should be the figure a reader
remembers: monotone, sign-changing at the calibration point, and larger than the
correction in item~3.

\item \textbf{Venous-pressure protocol: separation of the constitutive states.}
Sweep $P_{\rm RA}$ over the clinical range at fixed geometry and fixed driver,
for the four states. Report, per outlet and in total: mean coronary flow,
venular shear rate, $\mu_{\rm ap}$ ratio to healthy, and transmural pressure.
Expected outcome [venular $\dot{\gamma}$ from $400$ to
$\approx\!25\,\mathrm{s}^{-1}$ as $P_{\rm RA}$ goes from $4.5$ to
$14\,\mathrm{mmHg}$, with the hyperviscous/healthy viscosity ratio rising from
$1.17$ to $1.38$]. Two panels: shear rate and flow against $P_{\rm RA}$.
\emph{State explicitly the $\dot{\gamma}$ below which the equilibrium assumption
of \S\ref{sec:calibration-point} no longer holds, and mark it on the axis.}

\item \textbf{Waveform response.} For the healthy and hyperviscous states at
low and elevated $P_{\rm RA}$: coronary inflow $Q_i(t)$, venous outflow
$Q_d^{(i)}(t)$ and $p_{tm}^{(i)}(t)$ over one cycle. Metrics: systolic fraction
of the inflow volume, diastolic-to-systolic peak ratio, peak-to-mean venous
outflow, and the per-beat volume displaced by the intramyocardial compartment.
Verify $p_{tm}>0$ throughout, which certifies the non-negativity of the stored
volume without a penalty.

\item \textbf{Constitutive-model comparison at fixed state.} The same protocol
with power-law, Cross, Carreau--Yasuda and Quemada for the healthy fit. Expected
outcome: Cross and Carreau--Yasuda nearly indistinguishable at the operating
point; the power-law progressively wrong at low shear, where it lacks a plateau
and diverges; Quemada distinguished by admitting haematocrit and plasma viscosity
as explicit parameters. \emph{This is where the case for Quemada belongs: not a
better fit, but the only one of the four that separates the aggregation axis from
the haematocrit axis, and the haematocrit axis is the one that moves the
high-shear plateau where the bed operates.}

\item \textbf{Numerical properties of the coupling.} Amplification factor of the
lagged versus implicitly assembled outlet resistance, as a function of $\Delta t$
and $\mathcal{R}$; residual high-frequency content of the outlet flux in each
case; cost per accepted time step of the tabulated versus directly integrated
closure. Short, one table, but it is what makes the closure usable inside a 3D
solver and it answers the obvious referee question about stability.

\item \textbf{Control.} Verification that $p_{\rm ar}(t)$, $p_{\rm LV}(t)$ and
the ventricular volume are identical across all cases to machine precision, so
that every reported difference is attributable to the outlet closure. Two lines
of text, no figure. Cheap, and it forecloses an entire class of objection.

\end{enumerate}

\cbp{NOT to be reported, and why: a normotensive/hypertensive sweep of the
systemic circuit, and sensitivity to systemic vessel radius, length and wall
stiffness. Both probe the driver rather than the closure, and the previous draft
promised them; see \S\ref{sec:ApplicationI}.}

% ----------------------------------------------------------------------
% SECTION: SICKLE CELL ANEMIA (HbSS)
%


\begin{table}
    \begin{center}
    \begin{tabular}{l l c c c c c c c c}
        \hline
        Pathology & Model & \multicolumn{8}{c}{Parameters} \\
        \cline{3-10}
        & & $m$ & $n$ & $\mu_0$ & $\mu_\infty$ & $a$ & $\lambda$ & $\dot{\gamma}_c$ & $\mu_F$\\
        \hline
        Normal
         & PL & $5.0692$ & $0.6655$ & -- & -- & -- & -- & -- & -- \\
         & Q  & -- & -- & $80.3102$ & $5.0041$ & -- & -- & $2.2922$ & $1.7963$ \\
         & Cr & -- & $0.9503$ & $78.9514$ & $5.0629$ & -- & $4.1199$ & -- & -- \\
         & CY & -- & $0.2081$ & $71.9637$ & $4.8612$ & $1.4173$ & $6.1529$ & -- & -- \\
        [5pt] % Espacio vertical para separar categorías
        Hypertension
         & PL & $5.5256$ & $0.7048$ & -- & -- & -- & -- & -- & -- \\
         & Q  & -- & -- & $96.6447$ & $3.2835$ & -- & -- & $0.9992$ & $1.4502$ \\
         & Cr & -- & $0.8260$ & $92.5012$ & $4.1099$ & -- & $5.3469$ & -- & -- \\
         & CY & -- & $0.1703$ & $81.5805$ & $4.1796$ & $1.0324$ & $4.9438$ & -- & -- \\
        [5pt]
        Diabetes
         & PL & $3.9686$ & $0.6725$ & -- & -- & -- & -- & -- & -- \\
         & Q  & -- & -- & $122.2013$ & $4.7551$ & -- & -- & $1.7278$ & $1.9115$ \\
         & Cr & -- & $0.8633$ & $103.5405$ & $5.3610$ & -- & $10.4591$ & -- & -- \\
         & CY & -- & $0.1502$ & $104.6692$ & $5.3610$ & $0.8625$ & $11.1048$ & -- & -- \\
        [5pt]
        Hyperviscosity
         & PL & $8.2781$ & $0.5288$ & -- & -- & -- & -- & -- & -- \\
         & Q  & -- & -- & $313.0808$ & $2.6863$ & -- & -- & $14.9623$ & $1.3228$ \\
         & Cr & -- & $0.7982$ & $295.7832$ & $5.0291$ & -- & $13.7952$ & -- & -- \\
         & CY & -- & $0.2152$ & $319.4782$ & $5.5066$ & $0.7709$ & $16.9213$ & -- & -- \\
        \hline
    \end{tabular}
    \caption{Rheological parameters according to pathology and model.}
    \label{tab:parametros_reologicos}
    \end{center}
\end{table} 


\section{Conclusions}
\label{sec:conclusions}

\cbp{
WE NEED TO WRITE AT THE END. Outline of the claims the results should support,
in descending order of strength:

\begin{enumerate}
\item A constant-viscosity Windkessel resistance is exact only at the flow rate at
which it was calibrated, and its error grows monotonically away from that point,
reaching tens of per cent across the physiological range. The closure removes this
error and, with it, the unjustified choice of the constant.
\item The rheological operating point of an unresolved bed is not a modelling
parameter: \eqref{eq:gamma-identity} fixes it from calibre, effective path length
and pressure share, with the flow rate cancelling. This provides an a priori
criterion for deciding whether a non-Newtonian closure can matter in a given bed.
\item Applied to a coronary bed, that criterion predicts, and the simulations
confirm, that shear thinning contributes only a few per cent at rest and that a
hyperviscous constitutive state is weakly distinguishable from a healthy one. This
is reported as a result, not as a shortcoming.
\item The states do separate once the operating point is displaced, and elevated
venous drainage pressure is a physiologically grounded, geometry-preserving way to
do so.
\item The closure admits an exact derivative and a universal tabulation, which is
what makes it usable as a boundary condition inside a three-dimensional solver
without inner nonlinear solves.
\end{enumerate}

Limitations to state plainly rather than defend: the equilibrium assumption of the
purely viscous class fails for transients shorter than the aggregation time;
Fåhræus--Lindqvist effects are not contained in a bulk constitutive law applied to
vessels of tens of micrometres; a single effective segment per limb is a coarse
representation of a branching tree, as the transit-time check quantifies; and the
absence of vasomotor tone means the results describe a fixed-tone preparation, so
they should be compared against maximally dilated rather than autoregulated
measurements.
}

%%%%%%%%%%%%%%%%%%%%%%%%%%%%%%%%%%%%%%%%%%%%%%%





\appendix 
\section{Analytical evaluation of the flow-rate integral for the Cross and
Carreau--Yasuda models}\label{appA}

\subsection{Cross model}

Consider the flow-rate integral associated with the Cross model defined in
\eqref{Int:Cross}. Writing the integrand of the wall integral as
$\dot\gamma^{3}\,\mu(\dot\gamma)^{2}\,\mathrm d\tau/\mathrm d\dot\gamma$ and
expanding it algebraically, it decomposes into the sum of seven independent
contributions,
\begin{align}
I_1 &= \int_0^{\dot{\gamma}_w} \mu_{\infty}^3 \dot{\gamma}^3 \, d\dot{\gamma}
      = \frac{\mu_{\infty}^3 \dot{\gamma}_w^4}{4} \\[6pt]
I_2 &= 3\mu_{\infty}^2 \delta \int_0^{\dot{\gamma}_w}
      \frac{\dot{\gamma}^3}{1+(\lambda \dot{\gamma})^n} \, d\dot{\gamma} \\[6pt]
I_3 &= 3\mu_{\infty} \delta^2 \int_0^{\dot{\gamma}_w}
      \frac{\dot{\gamma}^3}{(1+(\lambda \dot{\gamma})^n)^2} \, d\dot{\gamma} \\[6pt]
I_4 &= -\mu_{\infty}^2 n \delta \int_0^{\dot{\gamma}_w}
      \frac{(\lambda \dot{\gamma})^{n}\dot{\gamma}^{3}}{(1+(\lambda \dot{\gamma})^n)^2}
      \, d\dot{\gamma} \\[6pt]
I_5 &= \delta^3 \int_0^{\dot{\gamma}_w}
      \frac{\dot{\gamma}^3}{(1+(\lambda \dot{\gamma})^n)^3} \, d\dot{\gamma} \\[6pt]
I_6 &= -2\mu_{\infty} \delta^2 n \int_0^{\dot{\gamma}_w}
      \frac{(\lambda \dot{\gamma})^{n}\dot{\gamma}^{3}}{(1+(\lambda \dot{\gamma})^n)^3}
      \, d\dot{\gamma} \\[6pt]
I_7 &= -n \delta^3 \int_0^{\dot{\gamma}_w}
      \frac{(\lambda \dot{\gamma})^{n}\dot{\gamma}^{3}}{(1+(\lambda \dot{\gamma})^n)^4}
      \, d\dot{\gamma}
\end{align}
where $\delta=\mu_0-\mu_\infty$. The integrals $I_1$, $I_2$, $I_3$ and $I_5$
share a common structure and admit a closed form in terms of the Gauss
hypergeometric function,
\begin{equation}
J_k^{\mathrm{Cr}} = \int_0^{\dot{\gamma}_w}
\frac{\dot{\gamma}^3}{\bigl(1+(\lambda \dot{\gamma})^n\bigr)^k} \, d\dot{\gamma}
= \frac{\dot{\gamma}_w^4}{4}\,
{}_2F_1\!\left(k, \tfrac{4}{n}; \tfrac{4}{n}+1; -(\lambda \dot{\gamma}_w)^n\right).
\end{equation}
The integrals $I_4$, $I_6$ and $I_7$ carry an additional factor
$(\lambda\dot\gamma)^n$ in the numerator and can be grouped into a second family,
\begin{align}
K_k^{\mathrm{Cr}}
&= \int_0^{\dot{\gamma}_w}
   \frac{\dot{\gamma}^3 (\lambda \dot{\gamma})^n}{(1+(\lambda \dot{\gamma})^n)^k}
   \, d\dot{\gamma} \nonumber\\
&= \int_0^{\dot{\gamma}_w}
   \frac{\dot{\gamma}^3}{(1+(\lambda \dot{\gamma})^n)^{k-1}} \, d\dot{\gamma}
 - \int_0^{\dot{\gamma}_w}
   \frac{\dot{\gamma}^3}{(1+(\lambda \dot{\gamma})^n)^k} \, d\dot{\gamma}.
\end{align}
Using the algebraic identity
\[
\frac{(\lambda \dot{\gamma})^n}{\bigl(1+(\lambda \dot{\gamma})^n\bigr)^k}
= \frac{1}{\bigl(1+(\lambda \dot{\gamma})^n\bigr)^{k-1}}
- \frac{1}{\bigl(1+(\lambda \dot{\gamma})^n\bigr)^k},
\]
one obtains the recurrence relation
\begin{equation}
K_k^{\mathrm{Cr}} = J_{k-1}^{\mathrm{Cr}} - J_k^{\mathrm{Cr}}.
\end{equation}
Finally, rewriting every contribution in terms of the family $J_k^{\mathrm{Cr}}$,
the total flow-rate integral of the Cross model can be expressed as
\begin{align}
I_{Cr} &= \mu_{\infty}^3 J_0^{\mathrm{Cr}}
        + \bigl[\mu_{\infty}^2 \delta (3 - n)\bigr] J_1^{\mathrm{Cr}}
        + \bigl[\mu_{\infty} \delta^2 (3 - 2n) + \mu_{\infty}^2 n \delta\bigr]
          J_2^{\mathrm{Cr}} \nonumber\\
       &\quad + \bigl[\delta^3 (1 - n) + 2\mu_{\infty} \delta^2 n\bigr]
          J_3^{\mathrm{Cr}} + n \delta^3 \, J_4^{\mathrm{Cr}}.
\end{align}

\subsection{Carreau--Yasuda model}

For the Carreau--Yasuda model the evaluation follows the same strategy. The wall
integral is $I_{\mathrm{CY}}=\int_0^{\dot\gamma_w}\dot\gamma^{3}\,
\mu(\dot\gamma)^{2}\,(\mathrm d\tau/\mathrm d\dot\gamma)\,\mathrm d\dot\gamma$, and
with $\alpha = 1+(\lambda\dot{\gamma})^a$ and $\delta=\mu_0-\mu_\infty$ the
derivative of the shear stress $\tau=\mu(\dot\gamma)\dot\gamma$ is
\begin{equation}
\frac{\mathrm d\tau}{\mathrm d\dot\gamma}
= \mu_\infty + \delta\,\alpha^{\frac{n-1}{a}}
+ \delta(n-1)(\lambda\dot\gamma)^a\,\alpha^{\frac{n-1-a}{a}}.
\label{eq:dtau_CY}
\end{equation}
Starting from the expression defined in \eqref{eq:I_CY_start}, this gives
\begin{align}
I_{\mathrm{CY}}
=&
\int_{0}^{\dot{\gamma}_w}
\dot{\gamma}^3
\left[\mu_{\infty}+\delta\,\alpha^{\frac{n-1}{a}}\right]^2 \cdot \nonumber\\
& \left[\mu_{\infty}
+\delta\,\alpha^{\frac{n-1}{a}}
+\delta(n-1)(\lambda \dot{\gamma})^a\,\alpha^{\frac{n-1-a}{a}} \right]\,\mathrm{d}
\dot{\gamma}.
\end{align}
Expanding the integrand algebraically yields
\begin{align}
I_{\mathrm{CY}}
= &\int \dot{\gamma}^3 \Bigg[ \mu_\infty^3
  + 3\mu_\infty^2 \delta\, \alpha^{\frac{n-1}{a}}
  + \mu_\infty^2 \delta (n-1)(\lambda \dot{\gamma})^a\, \alpha^{\frac{n-1}{a}-1}
  + 3\mu_\infty \delta^2\, \alpha^{\frac{2(n-1)}{a}} \nonumber\\
& + 2\mu_\infty \delta^2 (n-1)\, \alpha^{\frac{2(n-1)}{a}-1} (\lambda \dot{\gamma})^a
  + \delta^3\, \alpha^{\frac{3(n-1)}{a}}
  + \delta^3(n-1)\, \alpha^{\frac{3(n-1)}{a}-1} (\lambda \dot{\gamma})^a \Bigg]
  \,d\dot{\gamma},
\end{align}
which decomposes into seven independent integrals,
\begin{align}
I_1 &= \mu_\infty^3 \int_0^{\dot{\gamma}_w} \dot{\gamma}^3\, d\dot{\gamma} \\
I_2 &= 3\mu_\infty^2 \delta \int_0^{\dot{\gamma}_w}
       \dot{\gamma}^3 \bigl(1 + (\lambda \dot{\gamma})^a\bigr)^{\frac{n-1}{a}}
       d\dot{\gamma}\\
I_3 &= \mu_\infty^2 \delta (n-1) \int_0^{\dot{\gamma}_w}
       \dot{\gamma}^3 (\lambda\dot{\gamma})^a
       \bigl(1 + (\lambda \dot{\gamma})^a\bigr)^{\frac{n-1-a}{a}}
       d\dot{\gamma}\\
I_4 &= 3\mu_\infty \delta^2 \int_0^{\dot{\gamma}_w}
       \dot{\gamma}^3 \bigl(1 + (\lambda \dot{\gamma})^a\bigr)^{\frac{2n-2}{a}}
       d\dot{\gamma}\\
I_5 &= 2\mu_\infty \delta^2 (n-1) \int_0^{\dot{\gamma}_w}
       \dot{\gamma}^3 (\lambda\dot{\gamma})^a
       \bigl(1 + (\lambda \dot{\gamma})^a\bigr)^{\frac{2n-2-a}{a}}
       d\dot{\gamma}\\
I_6 &= \delta^3 \int_0^{\dot{\gamma}_w}
       \dot{\gamma}^3 \bigl(1 + (\lambda \dot{\gamma})^a\bigr)^{\frac{3n-3}{a}}
       d\dot{\gamma}\\
I_7 &= \delta^3 (n-1) \int_0^{\dot{\gamma}_w}
       \dot{\gamma}^3 (\lambda\dot{\gamma})^a
       \bigl(1 + (\lambda \dot{\gamma})^a\bigr)^{\frac{3n-3-a}{a}}
       d\dot{\gamma}
\end{align}
As in the Cross case, the identity
$(\lambda\dot{\gamma})^a = \bigl(1+(\lambda\dot{\gamma})^a\bigr) - 1$ allows every
integral to be written in terms of a single fundamental family,
\begin{equation}
J_k^{\mathrm{CY}} = \int_0^{\dot{\gamma}_w}
\dot{\gamma}^3 \bigl(1 + (\lambda \dot{\gamma})^a\bigr)^k\, d\dot{\gamma},
\end{equation}
which admits the closed form
\begin{equation}
J_k^{\mathrm{CY}} = \frac{\dot{\gamma}_w^4}{4}\,
{}_2F_1\!\left(-k, \tfrac{4}{a}; 1 + \tfrac{4}{a}; -(\lambda \dot{\gamma}_w)^a\right).
\end{equation}
Finally, expressing all contributions in terms of $J_k^{\mathrm{CY}}$ (with the
shorthand $k=(n-1)/a$), the flow-rate integral of the Carreau--Yasuda model reads
\begin{align}
I_{\mathrm{CY}}= &\mu_\infty^3 J_0^{\mathrm{CY}}
   + \mu_\infty^2 \delta (n+2) J_k^{\mathrm{CY}}
   - \mu_\infty^2 \delta (n-1) J_{k-1}^{\mathrm{CY}}
   + \mu_\infty \delta^2 (2n+1) J_{2k}^{\mathrm{CY}} \nonumber\\
   & - 2\mu_\infty \delta^2 (n-1) J_{2k-1}^{\mathrm{CY}}
   + \delta^3 n J_{3k}^{\mathrm{CY}}
   - \delta^3 (n-1) J_{3k-1}^{\mathrm{CY}}.
\end{align}

\section{Parameters and changes of variable used in the flow-rate integral of the
Quemada model}\label{appB}

For a flow governed by the Quemada rheology, the flow rate is expressed through a
function $F(\alpha,q)$. The dimensionless parameters follow from the changes of
variable
\begin{equation}
\xi = \frac{r}{R};\qquad
\alpha = \frac{\sqrt{\tau_0} + \sqrt{\mu_\infty \lambda}}{\sqrt{\tau_w}}, \qquad
q = \frac{\sqrt{\tau_0} - \sqrt{\mu_\infty \lambda}}{\sqrt{\tau_0} + \sqrt{\mu_\infty \lambda}},
\end{equation}
where the parameters $\tau_0$, $\mu_\infty$ and $\lambda$ are defined so as to
express the Quemada model in terms of stresses. Denoting the haematocrit
by~$\varphi$,
\begin{align}
\tau_0 &= \mu_p \gamma_c\,
  \frac{\bigl[\tfrac{1}{2}\varphi(k_0 - k_\infty)\bigr]^2}
       {\bigl(1 - \tfrac{1}{2} k_\infty \varphi\bigr)^4},\\[4pt]
\mu_\infty &= \frac{\mu_p}{\bigl(1 - \tfrac{1}{2} k_\infty \varphi\bigr)^2},\\[4pt]
\lambda &= \gamma_c
  \left[\frac{1 - \tfrac{1}{2} k_0 \varphi}{1 - \tfrac{1}{2} k_\infty \varphi}\right]^2 .
\end{align}
The function $F(\alpha,q)$ is a polynomial description in the variable $q$, whose
coefficients are
\begin{align}
P_1 &= -\tfrac{1}{7}(q+8), \\
P_2 &= -\tfrac{1}{42}(13q^2 - 8q - 7), \\
P_3 &= -\tfrac{1}{210}(143q^3 - 88q^2 - 113q + 48), \\
P_4 &= -\tfrac{1}{840}(1287q^4 - 792q^3 - 1342q^2 + 632q + 175), \\
P_5 &= -\tfrac{1}{840}(3003q^5 - 1848q^4 - 3894q^3 + 1944q^2 + 1011q - 256), \\
P_6 &= -\tfrac{1}{1680}(15015q^6 - 9240q^5 - 23331q^4 + 12096q^3 + 9081q^2
       - 3179q - 525), \\
P_7 &= -\tfrac{1}{1680}(45045q^7 - 27720q^6 - 82005q^5 + 43680q^4 + 42819q^3
       - 17304q^2 \nonumber\\
    &\qquad\quad - 5619q + 1024), \\
P_8 &= -\tfrac{1}{16}(1-q)^2(1+q)(429q^5 + 165q^4 - 330q^3 - 90q^2 + 45q + 5).
\end{align}
\section{}\label{appC}


\bibliographystyle{jfm}
\bibliography{jfm}
%Use of the above commands will create a bibliography using the .bib file. Shown below is a bibliography built from individual items.

%\bibliographystyle{jfm}
%\bibliography{jfm2esam}

%\begin{thebibliography}{99}

% \expandafter\ifx\csname natexlab\endcsname\relax
% \def\natexlab#1{#1}\fi
% \expandafter\ifx\csname selectlanguage\endcsname\relax
% \def\selectlanguage#1{\relax}\fi

% \bibitem[Batchelor (1971)]{Batchelor59}
% {\sc Batchelor, G.K.} 1971 {Small-scale variation of convected quantities like temperature in turbulent fluid part1, general discussion and the case of small conductivity}, {\it J. Fluid Mech.}, {\bf 5}, pp. 3-113-133.

% \bibitem [Bouguet (2008)]{Bouguet01}
% {\sc Bouguet, J.-Y} 2008 Camera Calibration Toolbox for Matlab {\url{http://www.vision.caltech.edu/bouguetj/calib_doc/}}.

%  \bibitem[Briukhanovetal et al (1967)] {Briukhanovetal1967}
% {\sc Briukhanov, A. V.,   Grigorian, S. S., Miagkov,  S. M., Plam, M. Y.,   I. E. Shurova, I. E.,   Eglit, M. E. and Yakimov, Y. L.} 1967
% {On some new approaches to the dynamics of snow avalanches},
% {\it Physics of Snow and Ice,  Proceedings of the International Conference on Low Temperature Science}
% {Vol 1} pp. 1221--1241 {Institute of Low Temperature Science, Hokkaido University, Sapporo, Hokkaido, Japan}.

% \bibitem[Brownell (2004)]{Brownell04}
%  {\sc Brownell,  C.J.  and Su,  L.K.} 2004  {Planar measurements of differential diffusion in turbulent jets}, {\it AIAA Paper},  pp. 2004-2335.

% \bibitem[Brownell and Su (2007)] {Brownell07}
%   {\sc Brownell, C.J. and  Su, L.K.} 2007 {Scale relations and spatial spectra in a differentially diffusing jet}, {\it AIAA Paper}, pp 2007-1314.

% \bibitem [Dennis (1985)] {Dennis85}
%  {\sc  Dennis, S.C.R.} 1985 {Compact explicit finite difference approximations to the Navier--Stokes equation},  { In \it Ninth Intl Conf. on Numerical Methods in Fluid Dynamics},  {ed Soubbaramayer and J.P. Boujot},  {Vol 218}, {\it Lecture Notes in Physics}, pp. 23-51. Springer.

% \bibitem [Edwards et al. (2017)]{EdwardsVirouletKokelaarGray2017}
% {\sc Edwards, A. N., Viroulet, S., Kokelaar, B. P. and Gray, J. M. N. T.} 2017 Formation of levees, troughs and elevated channels by avalanches on erodible slopes {\it J. Fluid Mech.}, {\bf 823}, pp. 278-315.

% \bibitem[Hwang et al (1970)] {Hwang70}
%  {\sc Hwang,  L.-S.  and  Tuck, E.O.} 1970 On the oscillations of harbours of arbitrary shape {\it J.~Fluid Mech.}, {\bf42}, pp 447-464.

% \bibitem[Josep and Saut (1990)] {JosephSaut1990}
%  {\sc Joseph, Daniel D. and Saut, Jean Claude} 1990 Short-wave instabilities and ill-posed initial-value problems {\it Theoretical and Computational Fluid Dynamics}, {\bf 1},  pp.191--227,  {\url{http://dx.doi.org/10.1007/BF00418002}}.

% \bibitem[Worster (1992)] {Worster92}
% { \sc  Worster, M.G.} 1992 The dynamics of mushy layers {\it Interactive dynamics of convection and solidification},
% {(ed. S.H. Davis and H.E. Huppert and W. Muller and M.G. Worster)}, pp. 113--138 {Kluwer}.

% \bibitem[Koch(1983)] {Koch83}
% {\sc Koch, W.} 1983 Resonant acoustic frequencies of flat plate cascades {\it J.~Sound Vib.}, {\bf 88}, pp. 233-242.

% \bibitem[Lee(1971)] {Lee71}
% {\sc Lee,  J.-J.}  1971 Wave-induced oscillations in harbours of arbitrary geometry {\it J.~Fluid Mech.}, {\bf 45}, pp. 375-394.

% \bibitem[Linton and  Evans (1992)] {Linton92}
%  {\sc  Linton, C.M. and  Evans, D.V.} 1992 The radiation and scattering of surface waves by a vertical circular cylinder in a channel {\it Phil.\ Trans.\ R. Soc.\ Lond.}, {\bf 338}, pp. 325-357.

% \bibitem [Martin(1980] {Martin80}
%  {\sc  Martin, P.A.} 1980 On the null-field equations for the exterior problems of acoustics {\it Q.~J. Mech.\ Appl.\ Maths},{\bf 33}, pp. 385--396.

% \bibitem [Rogallo(1981)] {Rogallo81}
%  {\sc Rogallo,  R.S.} 1981 Numerical experiments in homogeneous turbulence  { {\it Tech. Rep.} 81835}  {NASA Tech.\ Mem}.

% \bibitem[Ursell(1950)] {Ursell50}
% {\sc  Ursell, F.} 1950 Surface waves on deep water in the presence of a submerged cylinder i {\it Proc.\ Camb.\ Phil.\ Soc.}, {\bf 46}, pp.141--152.

% \bibitem[Wijngaarden (1968)]{Wijngaarden68}
% {\sc van Wijngaarden, L.} 1968 On the oscillations Near and at resonance in open pipes {\it J.~Engng Maths},{\bf 2}, pp. 225--240.

% \bibitem[Miller (1991)]{Miller91}
% { \sc  Miller, P.L.} 1991 Mixing in high Schmidt number turbulent jets {school {PhD thesis}} {California Institute of Technology}.

% \end{thebibliography}

%% End of file `jfm2esam.bib'.


\end{document}
